\documentclass[../main.tex]{subfiles}
\begin{document}

\chapter{GPU Programming Introduction}
\lessoninfo{
  video={M3, lessons 1--4 (Introduction of CUDA Programming; Occupancy; Host Code; Stencil Operation with CUDA)},
  sources={\cite{harris2013sharedmem} (shared memory); NVIDIA Hopper whitepaper \cite{nvidia2022hopper} (thread block clusters)},
  textbook={K12 \cite{kim2012gpgpu} §1.3; A18 \cite{aamodt2018gpu} ch.~2; CPG \cite{nvidia2024cudaguide} §2 Programming Model, §3.2 CUDA Runtime, §4.2 Hardware Multithreading; BPG \cite{nvidia2026cudabestpractices} §11 Execution Configuration Optimizations; PMPP \cite{hwu2022pmpp} ch.~2--4, 7--8},
  keywords={kernel, host code, grid, thread block, thread index, execution configuration, shared memory, barrier synchronization, BSP, memory spaces, occupancy, compute capability, cudaMalloc, cudaMemcpy, stencil, tile, halo},
}

Lesson~1 introduced the GPU as a throughput machine built from streaming
multiprocessors that execute warps of threads, and Lesson~2 showed how
parallel work is expressed with OpenMP in shared memory and with MPI
across distributed memories.  This lesson turns both into concrete CUDA
programs.  It introduces the kernel and the hierarchy of threads, blocks
and grid in which it runs, the memory spaces those threads can see, and
the barrier with which the threads of a block synchronise; it then asks
how many blocks fit on one multiprocessor at once (occupancy), writes the
host code that moves data to and from the device, and ends with the
stencil of Lesson~2 rewritten for CUDA, where the boundary exchange of
MPI becomes the loading of a halo into shared memory.

\section{Introduction of CUDA programming}
\label{sec:l03-cuda}

A CUDA program is split into two parts that run on two processors: code
that runs on the CPU, and a function that the CPU hands to the GPU for
execution by many threads at once.\source{M3, CUDA programming
introduction; K12 §1.3; CPG §2.1 Kernels, §2.2 Thread Hierarchy, §2.3
Memory Hierarchy; PMPP ch.~2--3; A18 ch.~2.}

\subsection{Kernel and host code}

\begin{definition}[Kernel, host code, device code]
  A \term{kernel}[カーネル] is a function that runs on the GPU and is
  executed by many threads in parallel.  It is declared with the
  \texttt{\_\_global\_\_} qualifier.  The \term{host code}[ホストコード]
  is the part of the program that runs on the CPU (the \term*{host},
  Lesson~1): it sets up the data, manages the GPU's memory and launches
  the kernels.  The \term{device code}[デバイスコード] is the part that
  runs on the GPU, the \term{device}[デバイス]: the kernels and the
  functions they call.\margin{The qualifiers: \texttt{\_\_global\_\_} marks
  a kernel (called from the host, runs on the device);
  \texttt{\_\_device\_\_} a function called from device code;
  \texttt{\_\_host\_\_} an ordinary CPU function.  CPG §7.1 Function Execution Space Specifiers \cite{nvidia2024cudaguide}.}
\end{definition}

The first kernel adds two vectors.  Each thread computes one element: it
derives its own index \texttt{i} from built-in variables, checks that the
index is inside the array, and performs one addition.
\begin{lstlisting}[language=C++, numbers=none]
__global__ void vectorAdd(const float *a, const float *b,
                          float *c, int n) {
  int i = blockIdx.x * blockDim.x + threadIdx.x;
  if (i < n) c[i] = a[i] + b[i];
}
\end{lstlisting}
There is no loop over the elements: the loop of the serial program (the
loop parallelism of Lesson~2) has become the set of threads, one per
iteration.  The host code launches the kernel with the triple-angle-bracket
syntax, naming how many threads to create and how they are grouped:
\begin{lstlisting}[language=C++, numbers=none]
int threads = 256;
int blocks = (n + threads - 1) / threads;  // ceil(n / 256)
vectorAdd<<<blocks, threads>>>(d_a, d_b, d_c, n);
\end{lstlisting}

\begin{definition}[Kernel launch, execution configuration]
  A \term{kernel launch}[カーネル起動] is the statement
  \texttt{kernel<<<numBlocks, threadsPerBlock>>>(args)} by which the host
  starts a kernel on the device.  The part between the angle brackets is
  the \term{execution configuration}[実行構成]: the number of thread
  blocks and the number of threads per block.  The total number of
  threads created is their product.\tip{The number of blocks is the
  ceiling of $n$ over the block size, written in integer arithmetic as
  \texttt{(n + threads - 1) / threads}; the last block then has idle
  threads, which the bounds check \texttt{if (i < n)} keeps from writing
  past the array.}
\end{definition}

\subsection{SPMD and the thread index}

The threads of a kernel all run the same code: GPU programming is
\term*{SPMD} (Lesson~1).  A kernel that only prints a message, launched
with four threads, prints it four times, in no particular order.  To make
the threads do different work, each must know which one it is.
\begin{lstlisting}[language=C++, numbers=none]
__global__ void hello() {
  printf("hello %d\n", threadIdx.x);
}
// host:
hello<<<1, 4>>>();   // prints hello 0..3 in some order
\end{lstlisting}

\begin{definition}[Thread index]
  The \term{thread index}[スレッドインデックス] \texttt{threadIdx} is a
  built-in variable, different in every thread, that gives the thread's
  position within its block.  It has three components
  \texttt{threadIdx.x}, \texttt{.y} and \texttt{.z}, so that a block can
  be laid out as a line, a rectangle or a box of threads, matching the
  shape of the data (a vector, an image, a volume).
\end{definition}

\begin{figure}[tb]
  \centering
  % One kernel executed by four threads (SPMD): the kernel on the left, four
% threads, four output lines that differ only by threadIdx.x.
% Redraws the outline figures m3l1_spmd and m3l1_threadidx (merged).
% Usage: % One kernel executed by four threads (SPMD): the kernel on the left, four
% threads, four output lines that differ only by threadIdx.x.
% Redraws the outline figures m3l1_spmd and m3l1_threadidx (merged).
% Usage: % One kernel executed by four threads (SPMD): the kernel on the left, four
% threads, four output lines that differ only by threadIdx.x.
% Redraws the outline figures m3l1_spmd and m3l1_threadidx (merged).
% Usage: \input{figures/l03-spmd}   (body width, 120 mm; drawing ~85 mm)
\begin{tikzpicture}[x=1mm, y=1mm, font=\scriptsize\sffamily,
                    >={Stealth[length=1.4mm]}]
  \node[draw, fill=ln-box, align=left, anchor=west, inner sep=3pt,
        font=\scriptsize\ttfamily] (prog) at (0,0)
    {\_\_global\_\_ void hello() \{\\
     \ \ printf("hello \%d\textbackslash{}n",\\
     \ \ \ \ \ \ \ \ \ threadIdx.x);\\
     \}};
  \node[above=1mm of prog, text=ln-muted]
    {\iflangja{1つのカーネル}{one kernel}};
  \foreach \i/\y in {0/9, 1/3, 2/-3, 3/-9} {
    \draw[->, ln-accent, thick] (prog.east) -- (64,\y)
      node[right, draw, fill=white, inner sep=2pt, font=\scriptsize\ttfamily]
      {hello \i};
    \node[text=ln-muted, anchor=east] at (63,\y+2.6)
      {\iflangja{スレッド}{thread} \i};
  }
  \node[text=ln-muted, align=center, anchor=north] at (50,-13)
    {\iflangja{4スレッド，実行順序は不定}{4 threads, order unspecified}};
\end{tikzpicture}
   (body width, 120 mm; drawing ~85 mm)
\begin{tikzpicture}[x=1mm, y=1mm, font=\scriptsize\sffamily,
                    >={Stealth[length=1.4mm]}]
  \node[draw, fill=ln-box, align=left, anchor=west, inner sep=3pt,
        font=\scriptsize\ttfamily] (prog) at (0,0)
    {\_\_global\_\_ void hello() \{\\
     \ \ printf("hello \%d\textbackslash{}n",\\
     \ \ \ \ \ \ \ \ \ threadIdx.x);\\
     \}};
  \node[above=1mm of prog, text=ln-muted]
    {\iflangja{1つのカーネル}{one kernel}};
  \foreach \i/\y in {0/9, 1/3, 2/-3, 3/-9} {
    \draw[->, ln-accent, thick] (prog.east) -- (64,\y)
      node[right, draw, fill=white, inner sep=2pt, font=\scriptsize\ttfamily]
      {hello \i};
    \node[text=ln-muted, anchor=east] at (63,\y+2.6)
      {\iflangja{スレッド}{thread} \i};
  }
  \node[text=ln-muted, align=center, anchor=north] at (50,-13)
    {\iflangja{4スレッド，実行順序は不定}{4 threads, order unspecified}};
\end{tikzpicture}
   (body width, 120 mm; drawing ~85 mm)
\begin{tikzpicture}[x=1mm, y=1mm, font=\scriptsize\sffamily,
                    >={Stealth[length=1.4mm]}]
  \node[draw, fill=ln-box, align=left, anchor=west, inner sep=3pt,
        font=\scriptsize\ttfamily] (prog) at (0,0)
    {\_\_global\_\_ void hello() \{\\
     \ \ printf("hello \%d\textbackslash{}n",\\
     \ \ \ \ \ \ \ \ \ threadIdx.x);\\
     \}};
  \node[above=1mm of prog, text=ln-muted]
    {\iflangja{1つのカーネル}{one kernel}};
  \foreach \i/\y in {0/9, 1/3, 2/-3, 3/-9} {
    \draw[->, ln-accent, thick] (prog.east) -- (64,\y)
      node[right, draw, fill=white, inner sep=2pt, font=\scriptsize\ttfamily]
      {hello \i};
    \node[text=ln-muted, anchor=east] at (63,\y+2.6)
      {\iflangja{スレッド}{thread} \i};
  }
  \node[text=ln-muted, align=center, anchor=north] at (50,-13)
    {\iflangja{4スレッド，実行順序は不定}{4 threads, order unspecified}};
\end{tikzpicture}

  \caption{SPMD: one kernel, four threads.  Every thread runs the same
    code; only \texttt{threadIdx.x} tells them apart, and nothing fixes
    the order of their output.}
  \label{fig:l03-spmd}
\end{figure}

With the index, the simplest form of the vector addition is
\texttt{c[threadIdx.x] = a[threadIdx.x] + b[threadIdx.x]}: thread~0 adds
element~0, thread~1 element~1, and so on (\cref{fig:l03-spmd}).  This form
works only for a single block, which is why the kernel above uses two more
built-in variables.

\subsection{The thread hierarchy}

\begin{definition}[Thread block, grid]
  Threads are grouped into blocks and blocks into a grid.  A
  \term{thread block}[スレッドブロック] is a group of threads that are
  executed together on one streaming multiprocessor and can cooperate
  through shared memory and barriers; PTX calls it a
  \term{cooperative thread array (CTA)}[協調スレッド配列（CTA）]\index{CTA|see{cooperative thread array (CTA)}}.
  The \term{grid}[グリッド] is the set of all blocks of one kernel
  launch.  The \term{block index}[ブロックインデックス] \texttt{blockIdx}
  gives a block's position in the grid, and the built-in
  \texttt{blockDim} and \texttt{gridDim} hold the sizes of a block and of
  the grid in each dimension.  This three-level structure, thread, block,
  grid, is the \term{thread hierarchy}[スレッド階層]
  (\cref{fig:l03-thread-hierarchy}).\margin{Limits, CPG Table~21
  \cite{nvidia2024cudaguide}: 1024 threads per block; block $x, y \le
  1024$, $z \le 64$; grid $x \le 2^{31}-1$, $y, z \le 65535$ (CC
  5.0--9.0).}
\end{definition}

\begin{figure}[tb]
  \centering
  % The CUDA thread hierarchy: a grid of 3 x 2 thread blocks, each block of
% 4 x 2 threads; one block is zoomed on the right with its thread indices.
% Original (the course states the hierarchy in words; cf. CPG Figure 4).
% Usage: % The CUDA thread hierarchy: a grid of 3 x 2 thread blocks, each block of
% 4 x 2 threads; one block is zoomed on the right with its thread indices.
% Original (the course states the hierarchy in words; cf. CPG Figure 4).
% Usage: % The CUDA thread hierarchy: a grid of 3 x 2 thread blocks, each block of
% 4 x 2 threads; one block is zoomed on the right with its thread indices.
% Original (the course states the hierarchy in words; cf. CPG Figure 4).
% Usage: \input{figures/l03-thread-hierarchy}   (body width, 120 mm)
\begin{tikzpicture}[x=1mm, y=1mm, font=\scriptsize\sffamily,
                    >={Stealth[length=1.4mm]}]
  % grid of blocks
  \draw[ln-accent, thick, rounded corners=1pt] (0,0) rectangle (64,38);
  \node[anchor=south west, text=ln-accent, font=\scriptsize\sffamily\bfseries]
    at (0,38.5) {\iflangja{グリッド}{grid} \ \texttt{gridDim = (3,2)}};
  \foreach \bx in {0,1,2} \foreach \by in {0,1} {
    \draw[fill=ln-box] (3+\bx*20,3+\by*17) rectangle ++(18,14);
    \node[anchor=south west, text=ln-muted, inner sep=1pt]
      at (3+\bx*20,13.2+\by*17) {\iflangja{ブロック}{block} (\bx,\by)};
    \foreach \tx in {0,...,3} \foreach \ty in {0,1} {
      \fill[ln-accent!60] (4.5+\bx*20+\tx*4,4+\by*17+\ty*4) rectangle ++(3,3);
    }
  }
  % zoom lines
  \draw[ln-muted, densely dashed] (41,20) -- (74,38);
  \draw[ln-muted, densely dashed] (41,3) -- (74,0);
  \draw[ln-accent, thick] (23,3) rectangle (41,17);
  % one block zoomed
  \draw[ln-accent, thick, rounded corners=1pt] (74,0) rectangle (118,38);
  \node[anchor=south west, text=ln-accent, font=\scriptsize\sffamily\bfseries]
    at (74,38.5) {\iflangja{ブロック}{block} (1,0) \ \texttt{blockDim = (4,2)}};
  \foreach \tx in {0,...,3} \foreach \ty in {0,1} {
    \draw[fill=ln-accent!60, draw=white]
      (\iflangja{75.5+\tx*10.5}{77+\tx*10},5+\ty*14) rectangle ++(\iflangja{10.2}{9},12);
    \node[text=white, align=center] at (\iflangja{80.6+\tx*10.5}{81.5+\tx*10},11+\ty*14)
      {\iflangja{スレッド}{thread}\\(\tx,\ty)};
  }
  \node[anchor=north, text=ln-muted] at (96,4)
    {\texttt{threadIdx = (x,y)}};
\end{tikzpicture}
   (body width, 120 mm)
\begin{tikzpicture}[x=1mm, y=1mm, font=\scriptsize\sffamily,
                    >={Stealth[length=1.4mm]}]
  % grid of blocks
  \draw[ln-accent, thick, rounded corners=1pt] (0,0) rectangle (64,38);
  \node[anchor=south west, text=ln-accent, font=\scriptsize\sffamily\bfseries]
    at (0,38.5) {\iflangja{グリッド}{grid} \ \texttt{gridDim = (3,2)}};
  \foreach \bx in {0,1,2} \foreach \by in {0,1} {
    \draw[fill=ln-box] (3+\bx*20,3+\by*17) rectangle ++(18,14);
    \node[anchor=south west, text=ln-muted, inner sep=1pt]
      at (3+\bx*20,13.2+\by*17) {\iflangja{ブロック}{block} (\bx,\by)};
    \foreach \tx in {0,...,3} \foreach \ty in {0,1} {
      \fill[ln-accent!60] (4.5+\bx*20+\tx*4,4+\by*17+\ty*4) rectangle ++(3,3);
    }
  }
  % zoom lines
  \draw[ln-muted, densely dashed] (41,20) -- (74,38);
  \draw[ln-muted, densely dashed] (41,3) -- (74,0);
  \draw[ln-accent, thick] (23,3) rectangle (41,17);
  % one block zoomed
  \draw[ln-accent, thick, rounded corners=1pt] (74,0) rectangle (118,38);
  \node[anchor=south west, text=ln-accent, font=\scriptsize\sffamily\bfseries]
    at (74,38.5) {\iflangja{ブロック}{block} (1,0) \ \texttt{blockDim = (4,2)}};
  \foreach \tx in {0,...,3} \foreach \ty in {0,1} {
    \draw[fill=ln-accent!60, draw=white]
      (\iflangja{75.5+\tx*10.5}{77+\tx*10},5+\ty*14) rectangle ++(\iflangja{10.2}{9},12);
    \node[text=white, align=center] at (\iflangja{80.6+\tx*10.5}{81.5+\tx*10},11+\ty*14)
      {\iflangja{スレッド}{thread}\\(\tx,\ty)};
  }
  \node[anchor=north, text=ln-muted] at (96,4)
    {\texttt{threadIdx = (x,y)}};
\end{tikzpicture}
   (body width, 120 mm)
\begin{tikzpicture}[x=1mm, y=1mm, font=\scriptsize\sffamily,
                    >={Stealth[length=1.4mm]}]
  % grid of blocks
  \draw[ln-accent, thick, rounded corners=1pt] (0,0) rectangle (64,38);
  \node[anchor=south west, text=ln-accent, font=\scriptsize\sffamily\bfseries]
    at (0,38.5) {\iflangja{グリッド}{grid} \ \texttt{gridDim = (3,2)}};
  \foreach \bx in {0,1,2} \foreach \by in {0,1} {
    \draw[fill=ln-box] (3+\bx*20,3+\by*17) rectangle ++(18,14);
    \node[anchor=south west, text=ln-muted, inner sep=1pt]
      at (3+\bx*20,13.2+\by*17) {\iflangja{ブロック}{block} (\bx,\by)};
    \foreach \tx in {0,...,3} \foreach \ty in {0,1} {
      \fill[ln-accent!60] (4.5+\bx*20+\tx*4,4+\by*17+\ty*4) rectangle ++(3,3);
    }
  }
  % zoom lines
  \draw[ln-muted, densely dashed] (41,20) -- (74,38);
  \draw[ln-muted, densely dashed] (41,3) -- (74,0);
  \draw[ln-accent, thick] (23,3) rectangle (41,17);
  % one block zoomed
  \draw[ln-accent, thick, rounded corners=1pt] (74,0) rectangle (118,38);
  \node[anchor=south west, text=ln-accent, font=\scriptsize\sffamily\bfseries]
    at (74,38.5) {\iflangja{ブロック}{block} (1,0) \ \texttt{blockDim = (4,2)}};
  \foreach \tx in {0,...,3} \foreach \ty in {0,1} {
    \draw[fill=ln-accent!60, draw=white]
      (\iflangja{75.5+\tx*10.5}{77+\tx*10},5+\ty*14) rectangle ++(\iflangja{10.2}{9},12);
    \node[text=white, align=center] at (\iflangja{80.6+\tx*10.5}{81.5+\tx*10},11+\ty*14)
      {\iflangja{スレッド}{thread}\\(\tx,\ty)};
  }
  \node[anchor=north, text=ln-muted] at (96,4)
    {\texttt{threadIdx = (x,y)}};
\end{tikzpicture}

  \caption{The thread hierarchy.  A grid of $3\times2$ blocks, each of
    $4\times2$ threads; a thread is addressed by its \texttt{blockIdx}
    and its \texttt{threadIdx}.}
  \label{fig:l03-thread-hierarchy}
\end{figure}

The warp of Lesson~1 is absent from this hierarchy on purpose: a warp is a
microarchitectural unit, the group of 32 consecutive threads of a block
that the hardware issues together, and the programming model of early
CUDA hid it.\margin{Later CUDA versions expose the warp through
warp-level primitives (shuffle, vote, Lesson~5); the hardware still cuts
every block into warps of \texttt{warpSize} = 32 consecutive threads,
Lesson~4.}  The block is the unit the programmer reasons about, and the
grid is the unit of a launch.

\paragraph{Global index.}  To give every thread of a multi-block grid a
distinct element, the thread index is combined with the block index.  For
a one-dimensional grid and block,
\begin{equation}
  \texttt{idx} = \texttt{blockIdx.x} \times \texttt{blockDim.x} +
  \texttt{threadIdx.x}.
  \label{eq:l03-global-index}
\end{equation}
\Cref{fig:l03-global-index} shows 16 elements processed by four blocks of
four threads: block~2 holds elements 8 to~11, and its thread~1 takes
element~9.

\begin{figure}[tb]
  \centering
  % Global index of 16 elements processed by 4 blocks of 4 threads:
% idx = blockIdx.x * blockDim.x + threadIdx.x.
% Redraws the outline figure m3l1_data_index (and m3l1_threadidx_vectoradd).
% Usage: % Global index of 16 elements processed by 4 blocks of 4 threads:
% idx = blockIdx.x * blockDim.x + threadIdx.x.
% Redraws the outline figure m3l1_data_index (and m3l1_threadidx_vectoradd).
% Usage: % Global index of 16 elements processed by 4 blocks of 4 threads:
% idx = blockIdx.x * blockDim.x + threadIdx.x.
% Redraws the outline figure m3l1_data_index (and m3l1_threadidx_vectoradd).
% Usage: \input{figures/l03-global-index}   (body width, 120 mm)
\begin{tikzpicture}[x=1mm, y=1mm, font=\scriptsize\sffamily,
                    >={Stealth[length=1.4mm]}]
  \def\w{6}
  \foreach \b/\col in {0/orange!45, 1/yellow!60, 2/blue!30, 3/green!45} {
    \foreach \t in {0,...,3} {
      \pgfmathtruncatemacro{\i}{\b*4+\t}
      \fill[\col, draw=white, line width=0.4pt]
        ({(\i)*\w},0) rectangle ({(\i+1)*\w},6);
      \node at ({(\i+0.5)*\w},3) {\i};
      \node[text=ln-muted] at ({(\i+0.5)*\w},-3) {\t};
    }
    \draw[ln-muted, thick] ({\b*4*\w},0) rectangle ({(\b+1)*4*\w},6);
    \node[align=center, text=ln-muted, anchor=north]
      at ({(\b*4+2)*\w},-5.5)
      {\iflangja{ブロック}{block} \b\\\texttt{blockIdx.x = \b}};
  }
  \node[anchor=east, text=ln-muted] at (-1,3) {\texttt{idx}};
  \node[anchor=east, text=ln-muted] at (-1,-3) {\texttt{threadIdx.x}};
  \node[anchor=east, text=ln-muted] at (-1,9.5) {\iflangja{要素}{element}};
  \foreach \i in {0,...,15} {
    \draw[->, ln-accent] ({(\i+0.5)*\w},12) -- ({(\i+0.5)*\w},6.6);
  }
  \node[anchor=north, text=ln-accent] at (48,-15.5)
    {\texttt{idx = blockIdx.x * blockDim.x + threadIdx.x}, \texttt{blockDim.x = 4}};
\end{tikzpicture}
   (body width, 120 mm)
\begin{tikzpicture}[x=1mm, y=1mm, font=\scriptsize\sffamily,
                    >={Stealth[length=1.4mm]}]
  \def\w{6}
  \foreach \b/\col in {0/orange!45, 1/yellow!60, 2/blue!30, 3/green!45} {
    \foreach \t in {0,...,3} {
      \pgfmathtruncatemacro{\i}{\b*4+\t}
      \fill[\col, draw=white, line width=0.4pt]
        ({(\i)*\w},0) rectangle ({(\i+1)*\w},6);
      \node at ({(\i+0.5)*\w},3) {\i};
      \node[text=ln-muted] at ({(\i+0.5)*\w},-3) {\t};
    }
    \draw[ln-muted, thick] ({\b*4*\w},0) rectangle ({(\b+1)*4*\w},6);
    \node[align=center, text=ln-muted, anchor=north]
      at ({(\b*4+2)*\w},-5.5)
      {\iflangja{ブロック}{block} \b\\\texttt{blockIdx.x = \b}};
  }
  \node[anchor=east, text=ln-muted] at (-1,3) {\texttt{idx}};
  \node[anchor=east, text=ln-muted] at (-1,-3) {\texttt{threadIdx.x}};
  \node[anchor=east, text=ln-muted] at (-1,9.5) {\iflangja{要素}{element}};
  \foreach \i in {0,...,15} {
    \draw[->, ln-accent] ({(\i+0.5)*\w},12) -- ({(\i+0.5)*\w},6.6);
  }
  \node[anchor=north, text=ln-accent] at (48,-15.5)
    {\texttt{idx = blockIdx.x * blockDim.x + threadIdx.x}, \texttt{blockDim.x = 4}};
\end{tikzpicture}
   (body width, 120 mm)
\begin{tikzpicture}[x=1mm, y=1mm, font=\scriptsize\sffamily,
                    >={Stealth[length=1.4mm]}]
  \def\w{6}
  \foreach \b/\col in {0/orange!45, 1/yellow!60, 2/blue!30, 3/green!45} {
    \foreach \t in {0,...,3} {
      \pgfmathtruncatemacro{\i}{\b*4+\t}
      \fill[\col, draw=white, line width=0.4pt]
        ({(\i)*\w},0) rectangle ({(\i+1)*\w},6);
      \node at ({(\i+0.5)*\w},3) {\i};
      \node[text=ln-muted] at ({(\i+0.5)*\w},-3) {\t};
    }
    \draw[ln-muted, thick] ({\b*4*\w},0) rectangle ({(\b+1)*4*\w},6);
    \node[align=center, text=ln-muted, anchor=north]
      at ({(\b*4+2)*\w},-5.5)
      {\iflangja{ブロック}{block} \b\\\texttt{blockIdx.x = \b}};
  }
  \node[anchor=east, text=ln-muted] at (-1,3) {\texttt{idx}};
  \node[anchor=east, text=ln-muted] at (-1,-3) {\texttt{threadIdx.x}};
  \node[anchor=east, text=ln-muted] at (-1,9.5) {\iflangja{要素}{element}};
  \foreach \i in {0,...,15} {
    \draw[->, ln-accent] ({(\i+0.5)*\w},12) -- ({(\i+0.5)*\w},6.6);
  }
  \node[anchor=north, text=ln-accent] at (48,-15.5)
    {\texttt{idx = blockIdx.x * blockDim.x + threadIdx.x}, \texttt{blockDim.x = 4}};
\end{tikzpicture}

  \caption{Global index of 16 elements over four blocks of four threads.
    The block index selects the group of four, the thread index the
    element within it.}
  \label{fig:l03-global-index}
\end{figure}

\begin{example}[Index arithmetic]
  \label{ex:l03-index}
  An array of 1000 floats is processed with 128 threads per block.
  \begin{itemize}
    \item Number of blocks: $\lceil 1000/128 \rceil = 8$, creating
      $8 \times 128 = 1024$ threads; the last 24 threads fail the bounds
      check and do nothing.
    \item The thread with \texttt{blockIdx.x = 5} and
      \texttt{threadIdx.x = 77} handles element $5 \times 128 + 77 = 717$.
    \item Element 900 is handled by block $\lfloor 900/128 \rfloor = 7$,
      thread $900 - 7 \times 128 = 4$.
  \end{itemize}
  For a two-dimensional image of width $W$ the same arithmetic is applied
  per dimension: \texttt{x = blockIdx.x * blockDim.x + threadIdx.x},
  \texttt{y = blockIdx.y * blockDim.y + threadIdx.y}, and the element is
  \texttt{img[y * W + x]}.\margin{Within a block, CPG §2.2 numbers the
  threads of a 2D block of size $(D_x, D_y)$ as $x + y D_x$: consecutive
  \texttt{threadIdx.x} values are consecutive threads, and therefore in
  the same warp.  \cite{nvidia2024cudaguide}.}
\end{example}

\subsection{Execution order}

Nothing in the model orders the threads or the blocks.  Four threads that
print their index may print 0, 1, 2, 3 or 3, 2, 1, 0 or 0, 2, 3, 1; the
six blocks of \cref{fig:l03-block-order} may run one after another in
index order, one after another in any other order, or several at a time.
A correct kernel produces the same result under every one of these
orders.\margin{CPG §2.2: thread blocks \enquote{are required to execute
independently: It must be possible to execute them in any order, in
parallel or in series.}  \cite{nvidia2024cudaguide}.  This independence is
what lets a kernel run unchanged on a GPU with 2 or with 132 SMs.}

\begin{figure}[tb]
  \centering
  % Three legal execution orders of the six blocks of a grid: in index order,
% in an arbitrary order, and two at a time.
% Redraws the outline figure m3l1_execution_order_blocks.
% Usage: % Three legal execution orders of the six blocks of a grid: in index order,
% in an arbitrary order, and two at a time.
% Redraws the outline figure m3l1_execution_order_blocks.
% Usage: % Three legal execution orders of the six blocks of a grid: in index order,
% in an arbitrary order, and two at a time.
% Redraws the outline figure m3l1_execution_order_blocks.
% Usage: \input{figures/l03-block-order}   (body width, 120 mm)
\begin{tikzpicture}[x=1mm, y=1mm, font=\scriptsize\sffamily,
                    >={Stealth[length=1.4mm]},
                    blk/.style={draw, fill=ln-box, minimum width=13mm,
                                minimum height=4.5mm, inner sep=1pt}]
  % the grid
  \draw[ln-accent, thick, rounded corners=1pt] (0,14) rectangle (22,30);
  \node[anchor=south, text=ln-accent] at (11,30.3) {\iflangja{グリッド}{grid}};
  \foreach \bx in {0,1,2} \foreach \by in {0,1} {
    \draw[fill=ln-box] (1.5+\bx*6.6,16+\by*6.5) rectangle ++(6,5.5);
    \node[font=\scriptsize\sffamily] at (4.5+\bx*6.6,18.7+\by*6.5) {\bx,\by};
  }
  \draw[->, ln-muted] (23,22) -- (27,22);
  % time axis
  \draw[->, ln-muted] (30,-2) -- (118,-2) node[below left, text=ln-muted]
    {\iflangja{時間}{time}};
  % case 1
  \node[anchor=east, text=ln-muted, align=right] at (29,33)
    {\iflangja{順序 1}{case 1}};
  \foreach \k/\t in {0/{0,0}, 1/{1,0}, 2/{2,0}, 3/{0,1}, 4/{1,1}, 5/{2,1}} {
    \node[blk, anchor=west] at (30+\k*14.5,33) {\t};
  }
  % case 2
  \node[anchor=east, text=ln-muted, align=right] at (29,25)
    {\iflangja{順序 2}{case 2}};
  \foreach \k/\t in {0/{2,1}, 1/{1,1}, 2/{1,0}, 3/{0,0}, 4/{2,0}, 5/{0,1}} {
    \node[blk, anchor=west] at (30+\k*14.5,25) {\t};
  }
  % case 3: three at a time
  \node[anchor=east, text=ln-muted, align=right] at (29,10)
    {\iflangja{順序 3}{case 3}};
  \foreach \k/\t in {0/{0,0}, 1/{1,0}, 2/{2,0}} {
    \node[blk, anchor=west, minimum width=40mm] at (30,16-\k*5.5) {\t};
  }
  \foreach \k/\t in {0/{0,1}, 1/{1,1}, 2/{2,1}} {
    \node[blk, anchor=west, minimum width=40mm] at (74,16-\k*5.5) {\t};
  }
\end{tikzpicture}
   (body width, 120 mm)
\begin{tikzpicture}[x=1mm, y=1mm, font=\scriptsize\sffamily,
                    >={Stealth[length=1.4mm]},
                    blk/.style={draw, fill=ln-box, minimum width=13mm,
                                minimum height=4.5mm, inner sep=1pt}]
  % the grid
  \draw[ln-accent, thick, rounded corners=1pt] (0,14) rectangle (22,30);
  \node[anchor=south, text=ln-accent] at (11,30.3) {\iflangja{グリッド}{grid}};
  \foreach \bx in {0,1,2} \foreach \by in {0,1} {
    \draw[fill=ln-box] (1.5+\bx*6.6,16+\by*6.5) rectangle ++(6,5.5);
    \node[font=\scriptsize\sffamily] at (4.5+\bx*6.6,18.7+\by*6.5) {\bx,\by};
  }
  \draw[->, ln-muted] (23,22) -- (27,22);
  % time axis
  \draw[->, ln-muted] (30,-2) -- (118,-2) node[below left, text=ln-muted]
    {\iflangja{時間}{time}};
  % case 1
  \node[anchor=east, text=ln-muted, align=right] at (29,33)
    {\iflangja{順序 1}{case 1}};
  \foreach \k/\t in {0/{0,0}, 1/{1,0}, 2/{2,0}, 3/{0,1}, 4/{1,1}, 5/{2,1}} {
    \node[blk, anchor=west] at (30+\k*14.5,33) {\t};
  }
  % case 2
  \node[anchor=east, text=ln-muted, align=right] at (29,25)
    {\iflangja{順序 2}{case 2}};
  \foreach \k/\t in {0/{2,1}, 1/{1,1}, 2/{1,0}, 3/{0,0}, 4/{2,0}, 5/{0,1}} {
    \node[blk, anchor=west] at (30+\k*14.5,25) {\t};
  }
  % case 3: three at a time
  \node[anchor=east, text=ln-muted, align=right] at (29,10)
    {\iflangja{順序 3}{case 3}};
  \foreach \k/\t in {0/{0,0}, 1/{1,0}, 2/{2,0}} {
    \node[blk, anchor=west, minimum width=40mm] at (30,16-\k*5.5) {\t};
  }
  \foreach \k/\t in {0/{0,1}, 1/{1,1}, 2/{2,1}} {
    \node[blk, anchor=west, minimum width=40mm] at (74,16-\k*5.5) {\t};
  }
\end{tikzpicture}
   (body width, 120 mm)
\begin{tikzpicture}[x=1mm, y=1mm, font=\scriptsize\sffamily,
                    >={Stealth[length=1.4mm]},
                    blk/.style={draw, fill=ln-box, minimum width=13mm,
                                minimum height=4.5mm, inner sep=1pt}]
  % the grid
  \draw[ln-accent, thick, rounded corners=1pt] (0,14) rectangle (22,30);
  \node[anchor=south, text=ln-accent] at (11,30.3) {\iflangja{グリッド}{grid}};
  \foreach \bx in {0,1,2} \foreach \by in {0,1} {
    \draw[fill=ln-box] (1.5+\bx*6.6,16+\by*6.5) rectangle ++(6,5.5);
    \node[font=\scriptsize\sffamily] at (4.5+\bx*6.6,18.7+\by*6.5) {\bx,\by};
  }
  \draw[->, ln-muted] (23,22) -- (27,22);
  % time axis
  \draw[->, ln-muted] (30,-2) -- (118,-2) node[below left, text=ln-muted]
    {\iflangja{時間}{time}};
  % case 1
  \node[anchor=east, text=ln-muted, align=right] at (29,33)
    {\iflangja{順序 1}{case 1}};
  \foreach \k/\t in {0/{0,0}, 1/{1,0}, 2/{2,0}, 3/{0,1}, 4/{1,1}, 5/{2,1}} {
    \node[blk, anchor=west] at (30+\k*14.5,33) {\t};
  }
  % case 2
  \node[anchor=east, text=ln-muted, align=right] at (29,25)
    {\iflangja{順序 2}{case 2}};
  \foreach \k/\t in {0/{2,1}, 1/{1,1}, 2/{1,0}, 3/{0,0}, 4/{2,0}, 5/{0,1}} {
    \node[blk, anchor=west] at (30+\k*14.5,25) {\t};
  }
  % case 3: three at a time
  \node[anchor=east, text=ln-muted, align=right] at (29,10)
    {\iflangja{順序 3}{case 3}};
  \foreach \k/\t in {0/{0,0}, 1/{1,0}, 2/{2,0}} {
    \node[blk, anchor=west, minimum width=40mm] at (30,16-\k*5.5) {\t};
  }
  \foreach \k/\t in {0/{0,1}, 1/{1,1}, 2/{2,1}} {
    \node[blk, anchor=west, minimum width=40mm] at (74,16-\k*5.5) {\t};
  }
\end{tikzpicture}

  \caption{Three legal executions of a grid of six blocks: in index order,
    in an arbitrary order, and three at a time.  The hardware chooses,
    and the choice may differ from one run to the next.}
  \label{fig:l03-block-order}
\end{figure}

The simplest mental model, used in this section, is that each block runs
on one streaming multiprocessor and the blocks are independent of each
other.  \Cref{sec:l03-occupancy} refines the first half: an SM usually
holds several blocks at once.

\subsection{Shared memory}

Because the blocks of a grid are independent, data must be divided among
them, and the threads of a block need a place in which to combine their
work.

\begin{definition}[Shared memory]
  The \term{shared memory}[共有メモリ] of a multiprocessor is a small,
  fast memory on the chip that is visible to all threads of one block and
  to no other block.  It is a \term{scratchpad memory}[スクラッチパッドメモリ]:
  software-managed, so the kernel itself copies data into it and out of
  it, as opposed to a cache that the hardware fills.  A variable is placed
  in shared memory with the \texttt{\_\_shared\_\_} qualifier, and lives
  as long as the block.\margin{Latencies on Tesla and Fermi GPUs, K12
  Table~1.4 \cite{kim2012gpgpu}: shared memory 4--32 cycles, global memory
  hundreds of cycles.  \cite{harris2013sharedmem} gives roughly
  $100\times$ lower latency for shared than for uncached global memory.}
\end{definition}

In the kernel below each thread writes one sum into shared memory, and
then reads the sum written by another thread:
\begin{lstlisting}[language=C++, numbers=none]
__global__ void addReverse(const int *a, const int *b,
                           int *c) {
  __shared__ int s[256];          // one block of 256
  int i = threadIdx.x;
  s[i] = a[i] + b[i];
  __syncthreads();                // all writes done
  c[i] = s[255 - i];              // another thread's sum
}
\end{lstlisting}
The read \texttt{s[255 - i]} is correct only if thread $255-i$ has
already stored its value; since threads run in no fixed order, the
kernel must wait for all of them.\caution{Lesson~2 used \enquote{shared
memory} in the CPU sense, one address space for all threads.  CUDA shared
memory is per block and on chip; a value written into it by one block is
invisible to every other block, and it is gone when the block finishes.}

\subsection{Synchronisation}

\begin{definition}[Barrier synchronisation]
  \texttt{\_\_syncthreads()} performs \term{barrier synchronization}[バリア同期]
  over the threads of one block: every thread of the block waits at the
  call until all of them have reached it, after which all continue.  It
  also guarantees that the shared and global memory writes made before the
  barrier are visible to the threads of the block after
  it.\caution{Every thread of the block must reach the same
  \texttt{\_\_syncthreads()}; a barrier inside an \texttt{if} that only
  some threads enter is undefined (in practice a hang).  Put the barrier
  outside the condition.}
\end{definition}

A barrier spans one block only.  Threads of different blocks cannot
synchronise with \texttt{\_\_syncthreads()}; when a computation needs all
blocks to have finished a phase before any block begins the next, the
phases are written as separate kernels, since a kernel launch completes
only after all of its blocks have.\margin{Since CUDA~9 a grid-wide
barrier exists through cooperative groups and a cooperative launch, which
needs every block resident at once (Lesson~5).}

\begin{definition}[Bulk synchronous parallel]
  The usual shape of a kernel that uses shared memory is three phases
  separated by barriers (\cref{fig:l03-bsp}): load the block's data from
  global memory into shared memory; \texttt{\_\_syncthreads()}; compute on
  the shared copy; \texttt{\_\_syncthreads()}; store the results back to
  global memory.  This alternation of a parallel phase and a barrier is
  the \term{bulk synchronous parallel (BSP)}[バルク同期並列（BSP）]\index{BSP|see{bulk synchronous parallel (BSP)}}
  pattern; the MPI stencil of Lesson~2, which alternates boundary
  exchange and computation, is the same pattern across processes.
\end{definition}

\begin{figure}[tb]
  \centering
  % Bulk synchronous use of shared memory: eight threads of one block load
% global memory into shared memory, wait at a barrier, compute, wait again,
% and store.  Dashed segments are threads waiting at __syncthreads().
% Redraws the outline figure m3l1_usage_with_thread_sync.
% Usage: % Bulk synchronous use of shared memory: eight threads of one block load
% global memory into shared memory, wait at a barrier, compute, wait again,
% and store.  Dashed segments are threads waiting at __syncthreads().
% Redraws the outline figure m3l1_usage_with_thread_sync.
% Usage: % Bulk synchronous use of shared memory: eight threads of one block load
% global memory into shared memory, wait at a barrier, compute, wait again,
% and store.  Dashed segments are threads waiting at __syncthreads().
% Redraws the outline figure m3l1_usage_with_thread_sync.
% Usage: \input{figures/l03-bsp}   (body width, 120 mm)
\begin{tikzpicture}[x=1mm, y=1mm, font=\scriptsize\sffamily,
                    >={Stealth[length=1.4mm]}]
  % phases: load 2..32, barrier 34, compute 36..70, barrier 72, store 74..110
  \foreach \t/\a/\b/\c in {0/26/60/100, 1/31/66/104, 2/20/58/108,
                           3/29/69/98, 4/23/55/106, 5/32/63/102,
                           6/18/68/109, 7/27/61/97} {
    \pgfmathsetmacro{\y}{30-\t*4}
    \node[anchor=east, text=ln-muted] at (1,\y) {t\t};
    \draw[ln-accent, thick] (2,\y) -- (\a,\y);
    \draw[ln-muted, densely dashed] (\a,\y) -- (34,\y);
    \draw[ln-accent, thick] (36,\y) -- (\b,\y);
    \draw[ln-muted, densely dashed] (\b,\y) -- (72,\y);
    \draw[ln-accent, thick] (74,\y) -- (\c,\y);
  }
  \foreach \x in {34, 72} {
    \draw[ln-caution, line width=1.2pt] (\x,-1) -- (\x,33);
    \node[anchor=south, text=ln-caution] at (\x,33.3)
      {\texttt{\_\_syncthreads()}};
  }
  \node[anchor=south, text=ln-muted] at (\iflangja{10}{18},33.3)
    {\iflangja{\shortstack{グローバル\\→ 共有}}{global $\to$ shared}};
  \node[anchor=south, text=ln-muted] at (53,33.3)
    {\iflangja{計算}{compute}};
  \node[anchor=south, text=ln-muted] at (\iflangja{94}{92},33.3)
    {\iflangja{\shortstack{共有 →\\グローバル}}{shared $\to$ global}};
  \draw[->, ln-muted] (2,-3) -- (112,-3) node[below left, text=ln-muted]
    {\iflangja{時間}{time}};
\end{tikzpicture}
   (body width, 120 mm)
\begin{tikzpicture}[x=1mm, y=1mm, font=\scriptsize\sffamily,
                    >={Stealth[length=1.4mm]}]
  % phases: load 2..32, barrier 34, compute 36..70, barrier 72, store 74..110
  \foreach \t/\a/\b/\c in {0/26/60/100, 1/31/66/104, 2/20/58/108,
                           3/29/69/98, 4/23/55/106, 5/32/63/102,
                           6/18/68/109, 7/27/61/97} {
    \pgfmathsetmacro{\y}{30-\t*4}
    \node[anchor=east, text=ln-muted] at (1,\y) {t\t};
    \draw[ln-accent, thick] (2,\y) -- (\a,\y);
    \draw[ln-muted, densely dashed] (\a,\y) -- (34,\y);
    \draw[ln-accent, thick] (36,\y) -- (\b,\y);
    \draw[ln-muted, densely dashed] (\b,\y) -- (72,\y);
    \draw[ln-accent, thick] (74,\y) -- (\c,\y);
  }
  \foreach \x in {34, 72} {
    \draw[ln-caution, line width=1.2pt] (\x,-1) -- (\x,33);
    \node[anchor=south, text=ln-caution] at (\x,33.3)
      {\texttt{\_\_syncthreads()}};
  }
  \node[anchor=south, text=ln-muted] at (\iflangja{10}{18},33.3)
    {\iflangja{\shortstack{グローバル\\→ 共有}}{global $\to$ shared}};
  \node[anchor=south, text=ln-muted] at (53,33.3)
    {\iflangja{計算}{compute}};
  \node[anchor=south, text=ln-muted] at (\iflangja{94}{92},33.3)
    {\iflangja{\shortstack{共有 →\\グローバル}}{shared $\to$ global}};
  \draw[->, ln-muted] (2,-3) -- (112,-3) node[below left, text=ln-muted]
    {\iflangja{時間}{time}};
\end{tikzpicture}
   (body width, 120 mm)
\begin{tikzpicture}[x=1mm, y=1mm, font=\scriptsize\sffamily,
                    >={Stealth[length=1.4mm]}]
  % phases: load 2..32, barrier 34, compute 36..70, barrier 72, store 74..110
  \foreach \t/\a/\b/\c in {0/26/60/100, 1/31/66/104, 2/20/58/108,
                           3/29/69/98, 4/23/55/106, 5/32/63/102,
                           6/18/68/109, 7/27/61/97} {
    \pgfmathsetmacro{\y}{30-\t*4}
    \node[anchor=east, text=ln-muted] at (1,\y) {t\t};
    \draw[ln-accent, thick] (2,\y) -- (\a,\y);
    \draw[ln-muted, densely dashed] (\a,\y) -- (34,\y);
    \draw[ln-accent, thick] (36,\y) -- (\b,\y);
    \draw[ln-muted, densely dashed] (\b,\y) -- (72,\y);
    \draw[ln-accent, thick] (74,\y) -- (\c,\y);
  }
  \foreach \x in {34, 72} {
    \draw[ln-caution, line width=1.2pt] (\x,-1) -- (\x,33);
    \node[anchor=south, text=ln-caution] at (\x,33.3)
      {\texttt{\_\_syncthreads()}};
  }
  \node[anchor=south, text=ln-muted] at (\iflangja{10}{18},33.3)
    {\iflangja{\shortstack{グローバル\\→ 共有}}{global $\to$ shared}};
  \node[anchor=south, text=ln-muted] at (53,33.3)
    {\iflangja{計算}{compute}};
  \node[anchor=south, text=ln-muted] at (\iflangja{94}{92},33.3)
    {\iflangja{\shortstack{共有 →\\グローバル}}{shared $\to$ global}};
  \draw[->, ln-muted] (2,-3) -- (112,-3) node[below left, text=ln-muted]
    {\iflangja{時間}{time}};
\end{tikzpicture}

  \caption{Bulk synchronous use of shared memory by eight threads of one
    block.  Threads finish each phase at different times (solid) and wait
    at the barrier (dashed) until the slowest arrives.}
  \label{fig:l03-bsp}
\end{figure}

The second barrier is needed when the compute phase writes its results
into shared memory and the store phase has threads copy values produced by
other threads, as in \cref{fig:l03-bsp}, or when the buffer is reloaded
for another tile (the next load would overwrite values that slow threads
are still reading).  A kernel whose compute phase writes each thread's
result straight to global memory can omit it.

\subsection{Several kernels}

A program with several distinct tasks uses several kernels, each with its
own execution configuration.  A serial step is a kernel of one block and
one thread; a data-parallel step is a kernel of many:
\begin{lstlisting}[language=C++, numbers=none]
sortKernel<<<1, 1>>>(d_data, n);            // serial step
addKernel<<<blocks, threads>>>(d_data, d_out, n);
storeKernel<<<1, 1>>>(d_out, n);            // serial step
\end{lstlisting}
Kernels launched from the same host thread run one after another by
default, so the second sees the results of the first.\margin{Running
kernels concurrently, or overlapping a kernel with a copy, needs CUDA
streams, Lesson~5.  A one-thread kernel uses one lane of one SM and leaves
the rest of the GPU idle: serial steps belong on the host unless the data
is already on the device.}

\subsection{Memory spaces}

A thread can address five memory spaces, distinguished by who can see
them and where they are (\cref{tab:l03-memory-spaces},
\cref{fig:l03-memory-spaces}).

\begin{definition}[Local, global, constant, texture memory]
  A thread's \term{local memory}[ローカルメモリ] is private to it; it
  holds the thread's variables that do not fit in registers.  Shared
  memory is private to one block.  The \term{global memory}[グローバルメモリ]
  is visible to every thread of every block of a kernel, and persists
  across kernel launches: it is where the host places the input and reads
  the output.  The \term{constant memory}[コンスタントメモリ] is a small
  read-only space, visible to all threads and cached, for values that do
  not change during a kernel, such as a physical constant or a filter.
  The \term{texture memory}[テクスチャメモリ] is a read-only space with a
  cache optimised for two-dimensional locality.  The last two are
  inherited from graphics, where they hold the constants of a scene and
  its texture images, and are rarely needed in general-purpose
  code.\margin{Sizes (CPG Table~21 \cite{nvidia2024cudaguide}): constant
  memory \SI{64}{\kibi\byte} per device with an \SI{8}{\kibi\byte} cache
  per SM; local memory up to \SI{512}{\kibi\byte} per thread.  Both live
  in device memory: local memory is as slow as global memory unless it
  hits in the L1 cache.}
\end{definition}

\begin{table}[tb]
  \centering
  \small
  \caption{The CUDA memory spaces: who can see them, how long they live,
    where they are.}
  \label{tab:l03-memory-spaces}
  \begin{tabular}{>{\raggedright\arraybackslash}p{18mm}
                  >{\raggedright\arraybackslash}p{24mm}
                  >{\raggedright\arraybackslash}p{16mm}
                  >{\raggedright\arraybackslash}p{42mm}}
    \toprule
    Space & Scope & Lifetime & Location, note \\
    \midrule
    registers & one thread & thread & on chip, allocated by the compiler \\
    local & one thread & thread & device memory (cached); spilled variables \\
    shared & one block & block & on chip, \texttt{\_\_shared\_\_} \\
    \addlinespace
    global & all blocks of all kernels & application & device memory; \texttt{cudaMalloc} \\
    constant & all blocks & application & device memory, read-only, cached; \texttt{\_\_constant\_\_} \\
    texture & all blocks & application & device memory, read-only, 2D cache \\
    \bottomrule
  \end{tabular}
\end{table}

\begin{figure}[tb]
  \centering
  % CUDA memory spaces.  Left: a GPU of 16 SMs, each with its own shared
% memory, above one global memory (with the read-only constant and texture
% spaces).  Right: the scope of each space: thread, block, grid.
% Redraws the outline figures m3l1_sharedmemory and m3l1_memory_space.
% Usage: % CUDA memory spaces.  Left: a GPU of 16 SMs, each with its own shared
% memory, above one global memory (with the read-only constant and texture
% spaces).  Right: the scope of each space: thread, block, grid.
% Redraws the outline figures m3l1_sharedmemory and m3l1_memory_space.
% Usage: % CUDA memory spaces.  Left: a GPU of 16 SMs, each with its own shared
% memory, above one global memory (with the read-only constant and texture
% spaces).  Right: the scope of each space: thread, block, grid.
% Redraws the outline figures m3l1_sharedmemory and m3l1_memory_space.
% Usage: \input{figures/l03-memory-spaces}   (body width, 120 mm)
\begin{tikzpicture}[x=1mm, y=1mm, font=\scriptsize\sffamily,
                    >={Stealth[length=1.4mm]}]
  % ---- left: SMs with shared memory, global memory
  \node[font=\small\sffamily\bfseries, anchor=south west] at (0,39) {GPU};
  \foreach \i in {0,...,3} \foreach \j in {0,...,3} {
    \draw[fill=ln-box] (\i*11,\j*7+11) rectangle ++(10,6);
    \node[anchor=west, inner sep=1pt] at (\i*11+0.5,\j*7+14.5) {SM};
    \fill[green!60!black!60] (\i*11+5.5,\j*7+11.8) rectangle ++(4,3.6);
  }
  \draw[fill=ln-accent!15] (0,2) rectangle (43,8);
  \node at (21.5,5) {\iflangja{グローバルメモリ}{global memory}};
  \node[anchor=north, text=ln-muted, align=center] at (21.5,1.5)
    {\iflangja{定数・テクスチャ（読み出し専用）}{constant, texture (read-only)}};
  \draw[green!60!black, ->] (49,36) -- (43.5,33.5);
  \node[anchor=west, text=green!60!black, align=left] at (48,36.5)
    {\iflangja{共有メモリ\\（SM 内，\\ブロック単位）}{shared memory\\(on chip, per block)}};
  % ---- right: scopes
  \begin{scope}[shift={(70,0)}]
    \draw[rounded corners=2pt, fill=ln-box] (0,0) rectangle (48,40);
    \node[anchor=north west, text=ln-accent, font=\scriptsize\sffamily\bfseries]
      at (1,39.5) {\iflangja{\shortstack[l]{グリッド\\（カーネル）}}{grid (kernel)}};
    \node[anchor=north east, align=right, text=ln-muted] at (47,39.5)
      {\iflangja{グローバル，定数，\\テクスチャメモリ}{global, constant,\\texture memory}};
    \draw[rounded corners=2pt, fill=white] (2,2) rectangle (46,27);
    \node[anchor=north west, text=ln-accent, font=\scriptsize\sffamily\bfseries]
      at (3,26.5) {\iflangja{ブロック}{block}};
    \node[anchor=north east, align=right, text=ln-muted] at (45,26.5)
      {\iflangja{共有メモリ}{shared memory}};
    \draw[rounded corners=2pt, fill=ln-accent!15] (4,4) rectangle (44,16);
    \node[anchor=north west, text=ln-accent, font=\scriptsize\sffamily\bfseries]
      at (5,15.5) {\iflangja{スレッド}{thread}};
    \node[anchor=north east, align=right, text=ln-muted] at (43,15.5)
      {\iflangja{レジスタ，\\ローカルメモリ}{registers,\\local memory}};
  \end{scope}
\end{tikzpicture}
   (body width, 120 mm)
\begin{tikzpicture}[x=1mm, y=1mm, font=\scriptsize\sffamily,
                    >={Stealth[length=1.4mm]}]
  % ---- left: SMs with shared memory, global memory
  \node[font=\small\sffamily\bfseries, anchor=south west] at (0,39) {GPU};
  \foreach \i in {0,...,3} \foreach \j in {0,...,3} {
    \draw[fill=ln-box] (\i*11,\j*7+11) rectangle ++(10,6);
    \node[anchor=west, inner sep=1pt] at (\i*11+0.5,\j*7+14.5) {SM};
    \fill[green!60!black!60] (\i*11+5.5,\j*7+11.8) rectangle ++(4,3.6);
  }
  \draw[fill=ln-accent!15] (0,2) rectangle (43,8);
  \node at (21.5,5) {\iflangja{グローバルメモリ}{global memory}};
  \node[anchor=north, text=ln-muted, align=center] at (21.5,1.5)
    {\iflangja{定数・テクスチャ（読み出し専用）}{constant, texture (read-only)}};
  \draw[green!60!black, ->] (49,36) -- (43.5,33.5);
  \node[anchor=west, text=green!60!black, align=left] at (48,36.5)
    {\iflangja{共有メモリ\\（SM 内，\\ブロック単位）}{shared memory\\(on chip, per block)}};
  % ---- right: scopes
  \begin{scope}[shift={(70,0)}]
    \draw[rounded corners=2pt, fill=ln-box] (0,0) rectangle (48,40);
    \node[anchor=north west, text=ln-accent, font=\scriptsize\sffamily\bfseries]
      at (1,39.5) {\iflangja{\shortstack[l]{グリッド\\（カーネル）}}{grid (kernel)}};
    \node[anchor=north east, align=right, text=ln-muted] at (47,39.5)
      {\iflangja{グローバル，定数，\\テクスチャメモリ}{global, constant,\\texture memory}};
    \draw[rounded corners=2pt, fill=white] (2,2) rectangle (46,27);
    \node[anchor=north west, text=ln-accent, font=\scriptsize\sffamily\bfseries]
      at (3,26.5) {\iflangja{ブロック}{block}};
    \node[anchor=north east, align=right, text=ln-muted] at (45,26.5)
      {\iflangja{共有メモリ}{shared memory}};
    \draw[rounded corners=2pt, fill=ln-accent!15] (4,4) rectangle (44,16);
    \node[anchor=north west, text=ln-accent, font=\scriptsize\sffamily\bfseries]
      at (5,15.5) {\iflangja{スレッド}{thread}};
    \node[anchor=north east, align=right, text=ln-muted] at (43,15.5)
      {\iflangja{レジスタ，\\ローカルメモリ}{registers,\\local memory}};
  \end{scope}
\end{tikzpicture}
   (body width, 120 mm)
\begin{tikzpicture}[x=1mm, y=1mm, font=\scriptsize\sffamily,
                    >={Stealth[length=1.4mm]}]
  % ---- left: SMs with shared memory, global memory
  \node[font=\small\sffamily\bfseries, anchor=south west] at (0,39) {GPU};
  \foreach \i in {0,...,3} \foreach \j in {0,...,3} {
    \draw[fill=ln-box] (\i*11,\j*7+11) rectangle ++(10,6);
    \node[anchor=west, inner sep=1pt] at (\i*11+0.5,\j*7+14.5) {SM};
    \fill[green!60!black!60] (\i*11+5.5,\j*7+11.8) rectangle ++(4,3.6);
  }
  \draw[fill=ln-accent!15] (0,2) rectangle (43,8);
  \node at (21.5,5) {\iflangja{グローバルメモリ}{global memory}};
  \node[anchor=north, text=ln-muted, align=center] at (21.5,1.5)
    {\iflangja{定数・テクスチャ（読み出し専用）}{constant, texture (read-only)}};
  \draw[green!60!black, ->] (49,36) -- (43.5,33.5);
  \node[anchor=west, text=green!60!black, align=left] at (48,36.5)
    {\iflangja{共有メモリ\\（SM 内，\\ブロック単位）}{shared memory\\(on chip, per block)}};
  % ---- right: scopes
  \begin{scope}[shift={(70,0)}]
    \draw[rounded corners=2pt, fill=ln-box] (0,0) rectangle (48,40);
    \node[anchor=north west, text=ln-accent, font=\scriptsize\sffamily\bfseries]
      at (1,39.5) {\iflangja{\shortstack[l]{グリッド\\（カーネル）}}{grid (kernel)}};
    \node[anchor=north east, align=right, text=ln-muted] at (47,39.5)
      {\iflangja{グローバル，定数，\\テクスチャメモリ}{global, constant,\\texture memory}};
    \draw[rounded corners=2pt, fill=white] (2,2) rectangle (46,27);
    \node[anchor=north west, text=ln-accent, font=\scriptsize\sffamily\bfseries]
      at (3,26.5) {\iflangja{ブロック}{block}};
    \node[anchor=north east, align=right, text=ln-muted] at (45,26.5)
      {\iflangja{共有メモリ}{shared memory}};
    \draw[rounded corners=2pt, fill=ln-accent!15] (4,4) rectangle (44,16);
    \node[anchor=north west, text=ln-accent, font=\scriptsize\sffamily\bfseries]
      at (5,15.5) {\iflangja{スレッド}{thread}};
    \node[anchor=north east, align=right, text=ln-muted] at (43,15.5)
      {\iflangja{レジスタ，\\ローカルメモリ}{registers,\\local memory}};
  \end{scope}
\end{tikzpicture}

  \caption{Left: each SM has its own shared memory; one global memory
    serves all of them.  Right: the scope of each space follows the
    thread hierarchy.}
  \label{fig:l03-memory-spaces}
\end{figure}

The hierarchy limits how information can be shared: data in a block's
shared memory stays in that SM, and two blocks communicate only through
global memory, in general across a kernel boundary.  Since Hopper, a
\term{thread block cluster}[スレッドブロッククラスタ] groups several
blocks that are scheduled together and can read and write each other's
shared memory, a new level between block and grid.\sidenote{Thread block
clusters (compute capability 9.0): the blocks of a cluster are guaranteed
to be co-scheduled on one GPU processing cluster and access each other's
shared memory as \emph{distributed shared memory}.  Up to 8 blocks per
cluster is the portable size; H100 supports up to 16 (CPG §2.2.1
\cite{nvidia2024cudaguide}; \cite{nvidia2022hopper}).  Lesson~5 returns
to them.}

\begin{keypoint}
  A CUDA program is a host program plus kernels.  A kernel runs as a
  grid of blocks of threads; each thread finds its data from
  \texttt{blockIdx}, \texttt{blockDim} and \texttt{threadIdx}, blocks
  run in no fixed order and must be independent, and threads of one
  block share an on-chip shared memory and synchronise with
  \texttt{\_\_syncthreads()}.
\end{keypoint}

\section{Occupancy}
\label{sec:l03-occupancy}

A streaming multiprocessor usually runs more than one block at a time.
How many, and why it matters, is the question of
occupancy.\source{M3, occupancy; BPG §11.1 Occupancy, §11.3 Thread and
Block Heuristics, §11.4 Effects of Shared Memory; CPG §4.2 Hardware
Multithreading; PMPP ch.~4.}

\subsection{Three resources}

\begin{definition}[Resource constraint]
  Every block that is resident on an SM holds a share of three of the
  SM's resources for its whole lifetime: a number of thread slots, a
  number of registers, and an amount of shared memory.  Each is a
  \term{resource constraint}[資源制約]: the number of blocks that can be
  resident at once is the smallest number allowed by any of the three
  (and by a fixed maximum count of blocks per SM).  A
  \term{register}[レジスタ] here is one 32-bit entry of the SM's register
  file; the registers a thread needs are fixed by the compiler, and a
  kernel that needs many is said to have high
  \term{register pressure}[レジスタ圧力].\margin{CC 7.0, 8.0 and 9.0 SMs:
  \num{65536} registers (\SI{256}{\kibi\byte}), 2048 thread slots
  (64 warps), 32 blocks, \SIrange{96}{228}{\kibi\byte} of shared
  memory; other CCs have fewer slots and blocks, \cref{tab:l03-cc-limits}.}
\end{definition}

Who decides each quantity is summarised in
\cref{tab:l03-who-sets}.  Threads per block is the programmer's choice at
launch.  Registers per thread is the compiler's: unlike a CPU, whose ISA
fixes the number of architectural registers (32 in most RISC ISAs), the
GPU lets the compiler allocate as many as the kernel needs, up to a
maximum, and the hardware then fits as many threads as the register file
allows.  Shared memory per block follows from the declarations:
\texttt{\_\_shared\_\_ int s[1000]} is $4 \times 1000 = 4000$ bytes.

\begin{table}[tb]
  \centering
  \small
  \caption{Who sets the per-block resource demand, and how to read it.}
  \label{tab:l03-who-sets}
  \begin{tabular}{>{\raggedright\arraybackslash}p{28mm}
                  >{\raggedright\arraybackslash}p{24mm}
                  >{\raggedright\arraybackslash}p{50mm}}
    \toprule
    Quantity & Set by & Where it is visible \\
    \midrule
    threads per block & host, at launch & the execution configuration \\
    registers per thread & compiler & \texttt{nvcc --ptxas-options=-v}; capped by \texttt{\_\_launch\_bounds\_\_} or \texttt{-maxrregcount} \\
    shared memory per block & \texttt{\_\_shared\_\_} declarations (plus the optional third launch parameter) & same compiler report \\
    \bottomrule
  \end{tabular}
\end{table}

\begin{example}[Blocks per SM]
  \label{ex:l03-blocks-per-sm}
  A multiprocessor offers 1024 thread slots, \num{32768} registers and
  \SI{48}{\kibi\byte} of shared memory.  A kernel is launched with 64
  threads per block, the compiler reports 40 registers per thread, and
  each block declares \SI{6}{\kibi\byte} of shared memory.
  \begin{description}
    \item[threads] $1024 / 64 = 16$ blocks;
    \item[registers] $\num{32768} / (40 \times 64) = 12.8$, so 12
      blocks;
    \item[shared memory] $48 / 6 = 8$ blocks.
  \end{description}
  The SM can hold $\min(16, 12, 8) = 8$ blocks, limited by shared
  memory.  Halving the shared memory per block would raise the limit to
  12, after which registers bind; only then would the register count
  matter.
\end{example}

\caution{The source's example quotes the register file as
\enquote{64 kilobytes}; the number meant is 64\,K $=$ \num{65536} registers,
which at 4 bytes each is \SI{256}{\kibi\byte}.  Register counts in
occupancy calculations are always counts of 32-bit registers, not bytes.}

\subsection{Why occupancy matters}

More resident blocks mean more resident warps, and more warps give the
SM more independent instructions to choose from while others wait for
memory.  \Cref{fig:l03-occupancy-cases} compares two ways of filling four
SMs: one block of 8 threads each (32 threads) or three blocks of 4 threads
each (48 threads).  The second keeps more threads in flight even though
each block is smaller, and so, usually, achieves more of the SM's
throughput.\margin{Why more warps hide memory latency is the subject of
Lesson~4 (multithreaded architecture); Lesson~7 shows the exception,
where too many warps thrash the L1 cache.}

\begin{figure}[tb]
  \centering
  % Two ways to fill four SMs: (a) one block of 8 threads per SM, 32 threads
% in total; (b) three blocks of 4 threads per SM, 48 threads in total.
% Redraws the outline figure m3l2_occupancy (with our own numbers).
% Usage: % Two ways to fill four SMs: (a) one block of 8 threads per SM, 32 threads
% in total; (b) three blocks of 4 threads per SM, 48 threads in total.
% Redraws the outline figure m3l2_occupancy (with our own numbers).
% Usage: % Two ways to fill four SMs: (a) one block of 8 threads per SM, 32 threads
% in total; (b) three blocks of 4 threads per SM, 48 threads in total.
% Redraws the outline figure m3l2_occupancy (with our own numbers).
% Usage: \input{figures/l03-occupancy-cases}   (body width, 120 mm)
\begin{tikzpicture}[x=1mm, y=1mm, font=\scriptsize\sffamily]
  % (a)
  \foreach \s in {0,...,3} {
    \draw[fill=ln-box] (\s*14,0) rectangle ++(13,22);
    \node[anchor=south, text=ln-muted] at (\s*14+6.5,22.3) {SM \s};
    \draw[ln-accent, thick] (\s*14+1,1) rectangle ++(11,20);
    \foreach \t in {0,...,7} {
      \pgfmathtruncatemacro{\r}{int(\t/2)}
      \pgfmathtruncatemacro{\c}{mod(\t,2)}
      \fill[ln-accent!60] (\s*14+2.5+\c*4.5,2.5+\r*4.5) rectangle ++(3.5,3.5);
    }
  }
  \node[anchor=north, align=center] at (27,-1.5)
    {(a) \iflangja{SM あたり 1 ブロック × 8 スレッド}{one block of 8 threads per SM}\\
     $4 \times 1 \times 8 = 32$ \iflangja{スレッド}{threads}};
  % (b)
  \begin{scope}[shift={(64,0)}]
    \foreach \s in {0,...,3} {
      \draw[fill=ln-box] (\s*14,0) rectangle ++(13,22);
      \node[anchor=south, text=ln-muted] at (\s*14+6.5,22.3) {SM \s};
      \foreach \b in {0,1,2} {
        \draw[ln-accent, thick] (\s*14+1,1+\b*7) rectangle ++(11,6);
        \foreach \t in {0,...,3} {
          \fill[ln-accent!60] (\s*14+1.8+\t*2.4,2.3+\b*7) rectangle ++(1.9,3.4);
        }
      }
    }
    \node[anchor=north, align=center] at (27,-1.5)
      {(b) \iflangja{SM あたり 3 ブロック × 4 スレッド}{three blocks of 4 threads per SM}\\
       $4 \times 3 \times 4 = 48$ \iflangja{スレッド}{threads}};
  \end{scope}
\end{tikzpicture}
   (body width, 120 mm)
\begin{tikzpicture}[x=1mm, y=1mm, font=\scriptsize\sffamily]
  % (a)
  \foreach \s in {0,...,3} {
    \draw[fill=ln-box] (\s*14,0) rectangle ++(13,22);
    \node[anchor=south, text=ln-muted] at (\s*14+6.5,22.3) {SM \s};
    \draw[ln-accent, thick] (\s*14+1,1) rectangle ++(11,20);
    \foreach \t in {0,...,7} {
      \pgfmathtruncatemacro{\r}{int(\t/2)}
      \pgfmathtruncatemacro{\c}{mod(\t,2)}
      \fill[ln-accent!60] (\s*14+2.5+\c*4.5,2.5+\r*4.5) rectangle ++(3.5,3.5);
    }
  }
  \node[anchor=north, align=center] at (27,-1.5)
    {(a) \iflangja{SM あたり 1 ブロック × 8 スレッド}{one block of 8 threads per SM}\\
     $4 \times 1 \times 8 = 32$ \iflangja{スレッド}{threads}};
  % (b)
  \begin{scope}[shift={(64,0)}]
    \foreach \s in {0,...,3} {
      \draw[fill=ln-box] (\s*14,0) rectangle ++(13,22);
      \node[anchor=south, text=ln-muted] at (\s*14+6.5,22.3) {SM \s};
      \foreach \b in {0,1,2} {
        \draw[ln-accent, thick] (\s*14+1,1+\b*7) rectangle ++(11,6);
        \foreach \t in {0,...,3} {
          \fill[ln-accent!60] (\s*14+1.8+\t*2.4,2.3+\b*7) rectangle ++(1.9,3.4);
        }
      }
    }
    \node[anchor=north, align=center] at (27,-1.5)
      {(b) \iflangja{SM あたり 3 ブロック × 4 スレッド}{three blocks of 4 threads per SM}\\
       $4 \times 3 \times 4 = 48$ \iflangja{スレッド}{threads}};
  \end{scope}
\end{tikzpicture}
   (body width, 120 mm)
\begin{tikzpicture}[x=1mm, y=1mm, font=\scriptsize\sffamily]
  % (a)
  \foreach \s in {0,...,3} {
    \draw[fill=ln-box] (\s*14,0) rectangle ++(13,22);
    \node[anchor=south, text=ln-muted] at (\s*14+6.5,22.3) {SM \s};
    \draw[ln-accent, thick] (\s*14+1,1) rectangle ++(11,20);
    \foreach \t in {0,...,7} {
      \pgfmathtruncatemacro{\r}{int(\t/2)}
      \pgfmathtruncatemacro{\c}{mod(\t,2)}
      \fill[ln-accent!60] (\s*14+2.5+\c*4.5,2.5+\r*4.5) rectangle ++(3.5,3.5);
    }
  }
  \node[anchor=north, align=center] at (27,-1.5)
    {(a) \iflangja{SM あたり 1 ブロック × 8 スレッド}{one block of 8 threads per SM}\\
     $4 \times 1 \times 8 = 32$ \iflangja{スレッド}{threads}};
  % (b)
  \begin{scope}[shift={(64,0)}]
    \foreach \s in {0,...,3} {
      \draw[fill=ln-box] (\s*14,0) rectangle ++(13,22);
      \node[anchor=south, text=ln-muted] at (\s*14+6.5,22.3) {SM \s};
      \foreach \b in {0,1,2} {
        \draw[ln-accent, thick] (\s*14+1,1+\b*7) rectangle ++(11,6);
        \foreach \t in {0,...,3} {
          \fill[ln-accent!60] (\s*14+1.8+\t*2.4,2.3+\b*7) rectangle ++(1.9,3.4);
        }
      }
    }
    \node[anchor=north, align=center] at (27,-1.5)
      {(b) \iflangja{SM あたり 3 ブロック × 4 スレッド}{three blocks of 4 threads per SM}\\
       $4 \times 3 \times 4 = 48$ \iflangja{スレッド}{threads}};
  \end{scope}
\end{tikzpicture}

  \caption{Filling four SMs.  (a)~One block of 8 threads per SM; (b)~three
    blocks of 4 threads per SM.  The smaller blocks keep 48 threads
    resident instead of 32.}
  \label{fig:l03-occupancy-cases}
\end{figure}

\begin{definition}[Occupancy, active warp]
  An \term{active warp}[アクティブワープ] is a warp that is resident on
  a multiprocessor, with its registers and shared memory allocated, and
  therefore a candidate for issue whenever its next instruction is ready;
  a warp that is ready in the current cycle is called \emph{eligible}.  The
  \term{occupancy}[占有率] of a kernel on a multiprocessor is the ratio of
  the number of active warps to the maximum number of warps the
  multiprocessor can hold \cite[§11.1]{nvidia2026cudabestpractices}.
  With $T$ threads per block, $R$ registers per thread and $S$ bytes of
  shared memory per block, on an SM with the limits of
  \cref{tab:l03-cc-limits},
  \begin{align}
    \text{blocks/SM} &= \min\!\left(
      \left\lfloor \frac{\text{Threads/SM}}{T} \right\rfloor,\;
      \text{Blocks/SM},\;
      \left\lfloor \frac{\text{Regs/SM}}{R\,T} \right\rfloor,\;
      \left\lfloor \frac{\text{Smem/SM}}{S} \right\rfloor \right),
    \label{eq:l03-blocks}\\
    \text{occupancy} &= \frac{\text{blocks/SM} \times \lceil T/32 \rceil}
                             {\text{Warps/SM}}.
    \label{eq:l03-occupancy}
  \end{align}
  \margin{\Cref{eq:l03-occupancy} gives the \emph{theoretical}
  occupancy; the \emph{achieved} occupancy that Nsight Compute measures
  at run time is lower when blocks finish unevenly or the grid is too
  small to fill every SM.}
\end{definition}

\begin{table}[tb]
  \centering
  \small
  \setlength{\tabcolsep}{4pt}
  \caption{Per-SM limits by compute capability (CPG Table~21
    \cite{nvidia2024cudaguide}; K = 1024, KiB binary).  The example GPUs
    are representative, not exhaustive.}
  \label{tab:l03-cc-limits}
  \begin{tabular}{>{\raggedright\arraybackslash}p{33mm}rrrrrr}
    \toprule
    Compute capability & 7.0 & 7.5 & 8.0 & 8.6 & 8.9 & 9.0 \\
    (example GPU) & V100 & T4 & A100 & RTX 3090 & L40 & H100 \\
    \midrule
    max threads per SM & 2048 & 1024 & 2048 & 1536 & 1536 & 2048 \\
    max warps per SM & 64 & 32 & 64 & 48 & 48 & 64 \\
    max blocks per SM & 32 & 16 & 32 & 16 & 24 & 32 \\
    32-bit registers per SM & 64\,K & 64\,K & 64\,K & 64\,K & 64\,K & 64\,K \\
    max registers per thread & 255 & 255 & 255 & 255 & 255 & 255 \\
    max shared memory per SM (KiB) & 96 & 64 & 164 & 100 & 100 & 228 \\
    max shared memory per block (KiB) & 96 & 64 & 163 & 99 & 99 & 227 \\
    \bottomrule
  \end{tabular}
\end{table}

\begin{definition}[Compute capability]
  The \term{compute capability}[Compute Capability] is NVIDIA's version
  number for the features and limits of a GPU generation, written
  \emph{major.minor}: 7.x Volta and Turing, 8.x Ampere and Ada, 9.0
  Hopper.  The limits of \cref{tab:l03-cc-limits} are per compute
  capability, which is why an occupancy calculation must name the
  target: the x.0 parts (V100, A100, H100) hold 64 warps, CC 7.5, 8.6
  and 8.9 (e.g.\ T4, RTX 3090, L40) only 32 or 48.\margin{On CC 8.x and
  9.0 the shared memory is carved out of a unified L1/shared array, and
  the split is configurable per kernel (CPG, Compute Capability 8.x
  \cite{nvidia2024cudaguide}); the table gives the maximum.  The warp
  scheduler that consumes the warps: Lesson~4.}
\end{definition}

\Needspace{14\baselineskip}
\begin{example}[Occupancy on an A100]
  \label{ex:l03-occupancy-a100}
  A kernel is launched with $T = 256$ threads per block (8 warps), uses
  $R = 64$ registers per thread and $S = \SI{48}{\kibi\byte}$ of shared
  memory per block, on a compute capability 8.0 SM.
  \begin{description}
    \item[threads] $2048 / 256 = 8$ blocks;
    \item[blocks] 32;
    \item[registers] $\num{65536} / (64 \times 256) = 4$ blocks;
    \item[shared memory] $\lfloor 164 / 48 \rfloor = 3$ blocks.
  \end{description}
  Three blocks are resident, $3 \times 8 = 24$ warps, and the occupancy is
  $24 / 64 = \SI{37.5}{\percent}$.  Shared memory binds; cutting it to
  \SI{32}{\kibi\byte} per block would allow 4 blocks (then registers bind)
  and $32/64 = \SI{50}{\percent}$.  Full occupancy (8 blocks) would need
  at most 32 registers per thread and at most about \SI{19}{\kibi\byte}
  of shared memory per block ($164/8 = 20.5$ less the
  \SI{1}{\kibi\byte} reservation).\margin{Granularity,
  ignored in \cref{eq:l03-blocks}: registers are allocated per warp in
  units of 256, and on CC 8.x and 9.0 CUDA reserves \SI{1}{\kibi\byte}
  of shared memory per block ($\lfloor 164/49 \rfloor$ is still~3).
  BPG §11.1.1; CPG, Compute Capability 8.x
  \cite{nvidia2026cudabestpractices,nvidia2024cudaguide}.}
\end{example}

In practice the arithmetic is done by the occupancy calculator of Nsight
Compute or by the runtime's occupancy API,
\begin{lstlisting}[language=C++, numbers=none]
cudaOccupancyMaxActiveBlocksPerMultiprocessor(
    &blocks, kernel, T, dynSmem);
\end{lstlisting}
where \texttt{dynSmem} is the dynamic shared memory per block (the static
\texttt{\_\_shared\_\_} usage is taken from the kernel itself); both tools
apply \cref{eq:l03-blocks} with the exact granularities of the
target device (BPG §11.1.1).

Higher occupancy is not always faster.  Beyond some point extra warps add
nothing, and a kernel limited by instruction throughput or by its own
memory bandwidth gains nothing from them; but low occupancy always
weakens the SM's ability to hide memory latency
\cite[§11.1]{nvidia2026cudabestpractices}.\tip{Starting points from BPG
§11.3: threads per block a multiple of the warp size; at least 64 per
block, and then only with several blocks per SM; 128--256 as the first
thing to try; several smaller blocks rather than one large block when
latency matters.}

\begin{keypoint}
  Blocks per SM $=$ the minimum over threads, registers, shared memory
  (and the block cap); occupancy $=$ resident warps over the maximum.
  Threads per block is chosen at launch, registers per thread by the
  compiler, shared memory by the declarations, and the limits by the
  compute capability.
\end{keypoint}

\section{Host code}
\label{sec:l03-host}

Shared memory is visible within a block and global memory to all blocks,
but neither is visible to the CPU: global memory sits in the device's own
memory, behind an I/O interface, and the host manages it through explicit
calls.\source{M3, host code; CPG §3.2.2 Device Memory; K12 §1.3.6.}

\begin{definition}[Host memory, device memory]
  The \term{host memory}[ホストメモリ] is the CPU's main memory, where
  the host program's variables live.  The \term{device memory}[デバイスメモリ]
  is the GPU's own memory (HBM or GDDR), in which global, constant,
  texture and local memory are placed.  The two are separate address
  spaces joined by an I/O interface (\cref{fig:l03-host-device}): a
  pointer into one is meaningless on the other side, and data crosses
  only by an explicit copy.\margin{The interface is the bottleneck.  H100: HBM3 reported as
  about \SI{3}{\tera\byte\per\second} (preliminary) \cite{nvidia2022hopper}
  against \SI{64}{\giga\byte\per\second} per direction of PCIe Gen~5 x16,
  about $47\times$.  A100 \SI{40}{\gibi\byte}: HBM2 at
  \SI{1555}{\giga\byte\per\second} \cite{nvidia2020ampere} against about
  \SI{32}{\giga\byte\per\second} per direction of Gen~4, about $49\times$.}
\end{definition}

\begin{figure}[tb]
  \centering
  % Host and device memory.  The CPU and its memory on the left, the GPU with
% its SMs and global memory on the right, an I/O interface between them;
% cudaMalloc reserves global memory, cudaMemcpy moves data across.
% Redraws the outline figure m3l3_managing_global_mem.
% Usage: % Host and device memory.  The CPU and its memory on the left, the GPU with
% its SMs and global memory on the right, an I/O interface between them;
% cudaMalloc reserves global memory, cudaMemcpy moves data across.
% Redraws the outline figure m3l3_managing_global_mem.
% Usage: % Host and device memory.  The CPU and its memory on the left, the GPU with
% its SMs and global memory on the right, an I/O interface between them;
% cudaMalloc reserves global memory, cudaMemcpy moves data across.
% Redraws the outline figure m3l3_managing_global_mem.
% Usage: \input{figures/l03-host-device}   (body width, 120 mm)
\begin{tikzpicture}[x=1mm, y=1mm, font=\scriptsize\sffamily,
                    >={Stealth[length=1.4mm]},
                    blk/.style={draw, fill=ln-box, inner sep=2pt, align=center}]
  % host
  \draw[ln-accent, thick, rounded corners=1pt] (0,0) rectangle (38,34);
  \node[anchor=south west, text=ln-accent, font=\scriptsize\sffamily\bfseries]
    at (0,34.5) {\iflangja{ホスト}{host}};
  \node[blk, minimum width=30mm, minimum height=9mm] at (19,26) {CPU};
  \node[blk, minimum width=30mm, minimum height=11mm, fill=ln-accent!15]
    (hm) at (19,10) {\iflangja{ホストメモリ}{host memory}\\\texttt{h\_a, h\_b, h\_c}};
  % interface
  \node[blk, minimum width=18mm, minimum height=34mm, fill=white] (io) at (54,17)
    {\iflangja{I/O\\インタフェース\\（PCIe）\\[20mm]\mbox{}}{I/O\\interface\\(PCIe)}};
  % device
  \draw[ln-accent, thick, rounded corners=1pt] (70,0) rectangle (118,34);
  \node[anchor=south west, text=ln-accent, font=\scriptsize\sffamily\bfseries]
    at (70,34.5) {\iflangja{デバイス}{device}};
  \foreach \i in {0,...,3} \foreach \j in {0,1} {
    \draw[fill=ln-box] (72+\i*11.5,20+\j*6.5) rectangle ++(10.5,5.5)
      node[midway] {SM};
  }
  \node[blk, minimum width=44mm, minimum height=11mm, fill=ln-accent!15]
    (dm) at (94,10) {\iflangja{グローバルメモリ（デバイスメモリ）}{global memory (device memory)}\\\texttt{d\_a, d\_b, d\_c}};
  % arrows
  \draw[->, thick, ln-accent] (34.2,\iflangja{15}{13}) -- (io.west |- 0,\iflangja{15}{13}) -- (io.east |- 0,\iflangja{15}{13}) -- (72,\iflangja{15}{13});
  \node[anchor=south, text=ln-accent] at (54,\iflangja{15.5}{13.5})
    {\iflangja{\shortstack{\texttt{cudaMemcpy}\\（ホスト → デバイス）}}{\texttt{cudaMemcpy} host $\to$ device}};
  \draw[<-, thick, ln-accent] (34.2,\iflangja{7.5}{6}) -- (io.west |- 0,\iflangja{7.5}{6}) -- (io.east |- 0,\iflangja{7.5}{6}) -- (72,\iflangja{7.5}{6});
  \node[anchor=north, text=ln-accent] at (54,\iflangja{7}{5.5})
    {\iflangja{\shortstack{\texttt{cudaMemcpy}\\（デバイス → ホスト）}}{\texttt{cudaMemcpy} device $\to$ host}};
  \node[anchor=north, text=ln-muted] at (94,4) {\texttt{cudaMalloc} / \texttt{cudaFree}};
\end{tikzpicture}
   (body width, 120 mm)
\begin{tikzpicture}[x=1mm, y=1mm, font=\scriptsize\sffamily,
                    >={Stealth[length=1.4mm]},
                    blk/.style={draw, fill=ln-box, inner sep=2pt, align=center}]
  % host
  \draw[ln-accent, thick, rounded corners=1pt] (0,0) rectangle (38,34);
  \node[anchor=south west, text=ln-accent, font=\scriptsize\sffamily\bfseries]
    at (0,34.5) {\iflangja{ホスト}{host}};
  \node[blk, minimum width=30mm, minimum height=9mm] at (19,26) {CPU};
  \node[blk, minimum width=30mm, minimum height=11mm, fill=ln-accent!15]
    (hm) at (19,10) {\iflangja{ホストメモリ}{host memory}\\\texttt{h\_a, h\_b, h\_c}};
  % interface
  \node[blk, minimum width=18mm, minimum height=34mm, fill=white] (io) at (54,17)
    {\iflangja{I/O\\インタフェース\\（PCIe）\\[20mm]\mbox{}}{I/O\\interface\\(PCIe)}};
  % device
  \draw[ln-accent, thick, rounded corners=1pt] (70,0) rectangle (118,34);
  \node[anchor=south west, text=ln-accent, font=\scriptsize\sffamily\bfseries]
    at (70,34.5) {\iflangja{デバイス}{device}};
  \foreach \i in {0,...,3} \foreach \j in {0,1} {
    \draw[fill=ln-box] (72+\i*11.5,20+\j*6.5) rectangle ++(10.5,5.5)
      node[midway] {SM};
  }
  \node[blk, minimum width=44mm, minimum height=11mm, fill=ln-accent!15]
    (dm) at (94,10) {\iflangja{グローバルメモリ（デバイスメモリ）}{global memory (device memory)}\\\texttt{d\_a, d\_b, d\_c}};
  % arrows
  \draw[->, thick, ln-accent] (34.2,\iflangja{15}{13}) -- (io.west |- 0,\iflangja{15}{13}) -- (io.east |- 0,\iflangja{15}{13}) -- (72,\iflangja{15}{13});
  \node[anchor=south, text=ln-accent] at (54,\iflangja{15.5}{13.5})
    {\iflangja{\shortstack{\texttt{cudaMemcpy}\\（ホスト → デバイス）}}{\texttt{cudaMemcpy} host $\to$ device}};
  \draw[<-, thick, ln-accent] (34.2,\iflangja{7.5}{6}) -- (io.west |- 0,\iflangja{7.5}{6}) -- (io.east |- 0,\iflangja{7.5}{6}) -- (72,\iflangja{7.5}{6});
  \node[anchor=north, text=ln-accent] at (54,\iflangja{7}{5.5})
    {\iflangja{\shortstack{\texttt{cudaMemcpy}\\（デバイス → ホスト）}}{\texttt{cudaMemcpy} device $\to$ host}};
  \node[anchor=north, text=ln-muted] at (94,4) {\texttt{cudaMalloc} / \texttt{cudaFree}};
\end{tikzpicture}
   (body width, 120 mm)
\begin{tikzpicture}[x=1mm, y=1mm, font=\scriptsize\sffamily,
                    >={Stealth[length=1.4mm]},
                    blk/.style={draw, fill=ln-box, inner sep=2pt, align=center}]
  % host
  \draw[ln-accent, thick, rounded corners=1pt] (0,0) rectangle (38,34);
  \node[anchor=south west, text=ln-accent, font=\scriptsize\sffamily\bfseries]
    at (0,34.5) {\iflangja{ホスト}{host}};
  \node[blk, minimum width=30mm, minimum height=9mm] at (19,26) {CPU};
  \node[blk, minimum width=30mm, minimum height=11mm, fill=ln-accent!15]
    (hm) at (19,10) {\iflangja{ホストメモリ}{host memory}\\\texttt{h\_a, h\_b, h\_c}};
  % interface
  \node[blk, minimum width=18mm, minimum height=34mm, fill=white] (io) at (54,17)
    {\iflangja{I/O\\インタフェース\\（PCIe）\\[20mm]\mbox{}}{I/O\\interface\\(PCIe)}};
  % device
  \draw[ln-accent, thick, rounded corners=1pt] (70,0) rectangle (118,34);
  \node[anchor=south west, text=ln-accent, font=\scriptsize\sffamily\bfseries]
    at (70,34.5) {\iflangja{デバイス}{device}};
  \foreach \i in {0,...,3} \foreach \j in {0,1} {
    \draw[fill=ln-box] (72+\i*11.5,20+\j*6.5) rectangle ++(10.5,5.5)
      node[midway] {SM};
  }
  \node[blk, minimum width=44mm, minimum height=11mm, fill=ln-accent!15]
    (dm) at (94,10) {\iflangja{グローバルメモリ（デバイスメモリ）}{global memory (device memory)}\\\texttt{d\_a, d\_b, d\_c}};
  % arrows
  \draw[->, thick, ln-accent] (34.2,\iflangja{15}{13}) -- (io.west |- 0,\iflangja{15}{13}) -- (io.east |- 0,\iflangja{15}{13}) -- (72,\iflangja{15}{13});
  \node[anchor=south, text=ln-accent] at (54,\iflangja{15.5}{13.5})
    {\iflangja{\shortstack{\texttt{cudaMemcpy}\\（ホスト → デバイス）}}{\texttt{cudaMemcpy} host $\to$ device}};
  \draw[<-, thick, ln-accent] (34.2,\iflangja{7.5}{6}) -- (io.west |- 0,\iflangja{7.5}{6}) -- (io.east |- 0,\iflangja{7.5}{6}) -- (72,\iflangja{7.5}{6});
  \node[anchor=north, text=ln-accent] at (54,\iflangja{7}{5.5})
    {\iflangja{\shortstack{\texttt{cudaMemcpy}\\（デバイス → ホスト）}}{\texttt{cudaMemcpy} device $\to$ host}};
  \node[anchor=north, text=ln-muted] at (94,4) {\texttt{cudaMalloc} / \texttt{cudaFree}};
\end{tikzpicture}

  \caption{Host and device memory.  \texttt{cudaMalloc} reserves global
    memory; \texttt{cudaMemcpy} moves data across the I/O interface in
    either direction.}
  \label{fig:l03-host-device}
\end{figure}

Allocating device memory, moving data to and from it, and releasing it
are three calls of the CUDA runtime:
\begin{lstlisting}[language=C++, numbers=none]
cudaMalloc((void **)&ptr, size);   // reserve size bytes
cudaMemcpy(dst, src, size, dir);   // copy size bytes
cudaFree(ptr);                     // release
\end{lstlisting}

\begin{definition}[Memory allocation and transfer]
  On the device, \term{memory allocation}[メモリ確保] is done by
  \texttt{cudaMalloc}, which reserves \texttt{size} bytes of global
  memory and returns a device pointer in \texttt{ptr}; \texttt{cudaFree}
  releases it.  \texttt{cudaMemcpy} copies \texttt{size} bytes in the
  direction \texttt{dir}, host to device
  (\texttt{cudaMemcpyHostToDevice}), device to host
  (\texttt{cudaMemcpyDeviceToHost}) or device to device.  Every runtime
  call returns a \texttt{cudaError\_t}; \texttt{cudaSuccess} means it
  worked.  \texttt{cudaMemcpy} is synchronous with respect to the host:
  a device-to-host copy returns when the data has arrived, and a copy
  issued after a kernel launch first waits for the kernel to
  finish.\margin{Asynchronous copies from pinned host
  memory, overlapped with kernels: Lesson~5.}
\end{definition}

The complete vector addition, kernel and host code, is then as follows
(the kernel is the one of \cref{sec:l03-cuda}).  The three
memory-management steps, allocate, copy in, copy out, surround the
launch; the host also checks the launch for an error.\caution{A kernel
launch returns to the host immediately and reports no error of its own.
\texttt{cudaGetLastError()} catches configuration errors (too many
threads per block, too many resources for the SM) at launch; an error
inside the kernel surfaces at the next synchronising call, such as the
\texttt{cudaMemcpy} that reads the result.  Check both return values in
real code.}
\begin{lstlisting}[language=C++, numbers=none]
#include <cstdio>
#include <cstdlib>
#include <cuda_runtime.h>

__global__ void vectorAdd(const float *a, const float *b,
                          float *c, int n) {
  int i = blockIdx.x * blockDim.x + threadIdx.x;
  if (i < n) c[i] = a[i] + b[i];
}

int main() {
  const int n = 1 << 20;                 // 1,048,576
  const size_t size = n * sizeof(float);
  float *h_a = (float *)malloc(size);
  float *h_b = (float *)malloc(size);
  float *h_c = (float *)malloc(size);
  for (int i = 0; i < n; ++i) {
    h_a[i] = i; h_b[i] = 2.0f * i;
  }

  float *d_a, *d_b, *d_c;                // device pointers
  cudaMalloc((void **)&d_a, size);
  cudaMalloc((void **)&d_b, size);
  cudaMalloc((void **)&d_c, size);
  cudaMemcpy(d_a, h_a, size, cudaMemcpyHostToDevice);
  cudaMemcpy(d_b, h_b, size, cudaMemcpyHostToDevice);

  int threads = 256;
  int blocks = (n + threads - 1) / threads;   // 4096
  vectorAdd<<<blocks, threads>>>(d_a, d_b, d_c, n);
  cudaError_t err = cudaGetLastError();
  if (err != cudaSuccess)
    printf("launch: %s\n", cudaGetErrorString(err));

  cudaMemcpy(h_c, d_c, size, cudaMemcpyDeviceToHost);
  printf("%f\n", h_c[1000]);             // 3000.000000

  cudaFree(d_a); cudaFree(d_b); cudaFree(d_c);
  free(h_a); free(h_b); free(h_c);
  return 0;
}
\end{lstlisting}

\begin{example}[Cost of the copies]
  \label{ex:l03-copy-cost}
  With $n = 2^{20}$ floats, each array is \SI{4}{\mebi\byte}: \SI{8}{\mebi\byte}
  go to the device and \SI{4}{\mebi\byte} come back, \SI{12}{\mebi\byte} in
  all, for $2^{20}$ additions.  At a sustained
  \SI{25}{\giga\byte\per\second} over PCIe the copies take about
  \SI{0.5}{\milli\second}; the kernel itself, reading \SI{8}{\mebi\byte}
  and writing \SI{4}{\mebi\byte} from device memory at
  \SI{1.5}{\tera\byte\per\second}, needs under
  \SI{10}{\micro\second}.  The transfer dominates by a factor of about
  60: a vector addition alone is not worth moving to the GPU, and data
  should stay on the device across as many kernels as possible.
\end{example}

\begin{definition}[Unified memory]
  With \term{unified memory}[ユニファイドメモリ]
  (\texttt{cudaMallocManaged}) a single pointer is usable from both host
  and device: the runtime
  and the GPU's memory management move pages between host and device
  memory on demand, and the explicit \texttt{cudaMemcpy} calls disappear.
  The copies still happen, as page migrations triggered by page faults,
  and their cost is the subject of Lesson~7.\margin{Lesson~7 gives the
  numbers: a GPU page fault costs tens of microseconds and the driver
  handles faults in batches, so unified memory simplifies the code but
  does not remove the transfer.}
\end{definition}

\section{Stencil operation with CUDA}
\label{sec:l03-stencil}

The stencil of Lesson~2 replaces each element $C$ of a grid by the
average of its four neighbours, $C' = (N + W + E + S)/4$.  Written for
CUDA it brings together everything above: a grid of blocks that divide
the data, shared memory for the values each block reuses, a barrier
between loading and computing, and a halo at the block
edges.\source{M3, stencil with CUDA; PMPP ch.~7 (convolution) and ch.~8
(stencil); BPG §10.2.3 Shared Memory.}

\subsection{Decomposition and reuse}

The data decomposition is the one of \cref{fig:l03-stencil-tile}, left:
each block owns one region of the image and each of its threads computes
one output element, so all outputs are computed in parallel.  Every interior
input element is read by four threads, those of its four neighbours
(five, if the stencil includes the centre).

\begin{definition}[Data reuse, tile, tiling]
  The \term{data reuse}[データ再利用] of a computation is the property
  that one value loaded from memory serves several of its operations.  A \term{tile}[タイル] is the
  rectangular sub-array that one block works on, and
  \term{tiling}[タイリング] is the technique of loading a tile once from
  global memory into shared memory and serving all the reuse from there,
  so that each global element is read about once per block instead of
  once per use.  The reuse pattern of a stencil is known in advance, which
  is what makes a software-managed scratchpad the right
  cache.\margin{The same tiling turns a matrix multiplication from one
  global load per multiply into one per $\text{tile size}$ multiplies;
  Lesson~5 treats it, along with the bank conflicts of Lesson~4 that a
  badly laid out tile can cause.}
\end{definition}

Each thread therefore follows the BSP shape of \cref{fig:l03-bsp}: load
one element from global memory into the block's shared tile;
\texttt{\_\_syncthreads()}; compute its stencil from the tile; write the
result to global memory.

\subsection{Boundaries and the halo}

Shared memory is per block, but an element at the edge of a block's
region has neighbours in the regions of other blocks: in
\cref{fig:l03-stencil-tile}, $C$ at the corner of block (0,1) needs
neighbour~1 from block (0,0) and neighbour~3 from block (1,1).  Those
blocks' shared memories are unreachable, so each block loads, in addition
to its own region, the one-element border around it, the
\term*{halo} of Lesson~2, from global memory into its tile
(\cref{fig:l03-stencil-tile}, right).  This is the boundary exchange of
the MPI stencil without the messages: the halo is simply read from global
memory, which every block can see, and no synchronisation between blocks
is needed because the input array is read-only during the kernel and the
outputs go to a separate array.

\begin{figure}[tb]
  \centering
  % A stencil on a 10 x 8 image, one thread block per quadrant (left); the
% element C and its four neighbours straddle three blocks.  Right: the
% tile of block (0,1) (bottom-left) with its halo: cells owned by other
% blocks (grey) are loaded from global memory, cells outside the image are
% zero-filled.
% Redraws the outline figures m3l4_stencil and m3l4_stencil_boundaries.
% Usage: % A stencil on a 10 x 8 image, one thread block per quadrant (left); the
% element C and its four neighbours straddle three blocks.  Right: the
% tile of block (0,1) (bottom-left) with its halo: cells owned by other
% blocks (grey) are loaded from global memory, cells outside the image are
% zero-filled.
% Redraws the outline figures m3l4_stencil and m3l4_stencil_boundaries.
% Usage: % A stencil on a 10 x 8 image, one thread block per quadrant (left); the
% element C and its four neighbours straddle three blocks.  Right: the
% tile of block (0,1) (bottom-left) with its halo: cells owned by other
% blocks (grey) are loaded from global memory, cells outside the image are
% zero-filled.
% Redraws the outline figures m3l4_stencil and m3l4_stencil_boundaries.
% Usage: \input{figures/l03-stencil-tile}   (body width, 120 mm)
\begin{tikzpicture}[x=1mm, y=1mm, font=\scriptsize\sffamily,
                    >={Stealth[length=1.4mm]}]
  \def\s{3.4}  % cell size
  % ---- left: image split among four blocks (rows counted from the top)
  \foreach \r in {1,...,8} {
    \foreach \c in {1,...,10} {
      \ifnum\r<5
        \ifnum\c<6 \def\col{orange!45}\else \def\col{yellow!60}\fi
      \else
        \ifnum\c<6 \def\col{blue!30}\else \def\col{green!45}\fi
      \fi
      \fill[\col, draw=white, line width=0.3pt]
        ({(\c-1)*\s},{-\r*\s}) rectangle ({\c*\s},{-(\r-1)*\s});
    }
  }
  \draw[ln-muted] (0,0) rectangle (10*\s,-8*\s);
  \draw[ln-muted, thick] (5*\s,0) -- (5*\s,-8*\s);
  \draw[ln-muted, thick] (0,-4*\s) -- (10*\s,-4*\s);
  \foreach \c/\r/\t in {5/5/C, 5/4/1, 4/5/2, 6/5/3, 5/6/4} {
    \node[font=\scriptsize\sffamily\bfseries] at ({(\c-0.5)*\s},{-(\r-0.5)*\s}) {\t};
  }
  \node[text=ln-muted] at (2.5*\s,1.8) {\iflangja{ブロック}{block} (0,0)};
  \node[text=ln-muted] at (7.5*\s,1.8) {\iflangja{ブロック}{block} (1,0)};
  \node[anchor=north, text=ln-muted] at (2.5*\s,-8*\s-0.5) {\iflangja{ブロック}{block} (0,1)};
  \node[anchor=north, text=ln-muted] at (7.5*\s,-8*\s-0.5) {\iflangja{ブロック}{block} (1,1)};
  % ---- right: the tile of block (0,1) with its halo
  \begin{scope}[shift={(56,0)}]
    % outside-the-image cells: left column and bottom row (zero)
    \foreach \r in {4,...,9} {
      \draw[ln-muted, densely dotted, fill=white] ({-1*\s},{-\r*\s}) rectangle ({0},{-(\r-1)*\s});
      \node[text=ln-muted] at ({-0.5*\s},{-(\r-0.5)*\s}) {0};
    }
    \foreach \c in {1,...,6} {
      \draw[ln-muted, densely dotted, fill=white] ({(\c-1)*\s},{-9*\s}) rectangle ({\c*\s},{-8*\s});
      \node[text=ln-muted] at ({(\c-0.5)*\s},{-8.5*\s}) {0};
    }
    % halo cells inside the image: top row from block (0,0), right column from (1,1)
    \foreach \c in {1,...,5} {
      \fill[gray!45, draw=white, line width=0.3pt] ({(\c-1)*\s},{-4*\s}) rectangle ({\c*\s},{-3*\s});
    }
    \fill[gray!45, draw=white, line width=0.3pt] ({5*\s},{-4*\s}) rectangle ({6*\s},{-3*\s});
    \foreach \r in {5,...,8} {
      \fill[gray!45, draw=white, line width=0.3pt] ({5*\s},{-\r*\s}) rectangle ({6*\s},{-(\r-1)*\s});
      \foreach \c in {1,...,5} {
        \fill[blue!30, draw=white, line width=0.3pt] ({(\c-1)*\s},{-\r*\s}) rectangle ({\c*\s},{-(\r-1)*\s});
      }
    }
    \draw[ln-muted] (0,-4*\s) rectangle (5*\s,-8*\s);
    \draw[ln-accent, thick] ({-1*\s},-3*\s) rectangle (6*\s,-9*\s);
    \foreach \c/\r/\t in {5/5/C, 5/4/1, 4/5/2, 6/5/3, 5/6/4} {
      \node[font=\scriptsize\sffamily\bfseries] at ({(\c-0.5)*\s},{-(\r-0.5)*\s}) {\t};
    }
    \node[anchor=south, text=ln-accent] at (2.5*\s,-3*\s+0.5)
      {\iflangja{共有メモリのタイル（ハローを含む）}{shared-memory tile with halo}};
    \node[anchor=west, text=ln-muted, align=left] at (6.5*\s,-5.5*\s)
      {\iflangja{灰色：他ブロックの要素，\\グローバルメモリから読む\\点線：画像の外，0 を入れる}{grey: other blocks' cells,\\read from global memory\\dotted: outside the image, zero}};
  \end{scope}
\end{tikzpicture}
   (body width, 120 mm)
\begin{tikzpicture}[x=1mm, y=1mm, font=\scriptsize\sffamily,
                    >={Stealth[length=1.4mm]}]
  \def\s{3.4}  % cell size
  % ---- left: image split among four blocks (rows counted from the top)
  \foreach \r in {1,...,8} {
    \foreach \c in {1,...,10} {
      \ifnum\r<5
        \ifnum\c<6 \def\col{orange!45}\else \def\col{yellow!60}\fi
      \else
        \ifnum\c<6 \def\col{blue!30}\else \def\col{green!45}\fi
      \fi
      \fill[\col, draw=white, line width=0.3pt]
        ({(\c-1)*\s},{-\r*\s}) rectangle ({\c*\s},{-(\r-1)*\s});
    }
  }
  \draw[ln-muted] (0,0) rectangle (10*\s,-8*\s);
  \draw[ln-muted, thick] (5*\s,0) -- (5*\s,-8*\s);
  \draw[ln-muted, thick] (0,-4*\s) -- (10*\s,-4*\s);
  \foreach \c/\r/\t in {5/5/C, 5/4/1, 4/5/2, 6/5/3, 5/6/4} {
    \node[font=\scriptsize\sffamily\bfseries] at ({(\c-0.5)*\s},{-(\r-0.5)*\s}) {\t};
  }
  \node[text=ln-muted] at (2.5*\s,1.8) {\iflangja{ブロック}{block} (0,0)};
  \node[text=ln-muted] at (7.5*\s,1.8) {\iflangja{ブロック}{block} (1,0)};
  \node[anchor=north, text=ln-muted] at (2.5*\s,-8*\s-0.5) {\iflangja{ブロック}{block} (0,1)};
  \node[anchor=north, text=ln-muted] at (7.5*\s,-8*\s-0.5) {\iflangja{ブロック}{block} (1,1)};
  % ---- right: the tile of block (0,1) with its halo
  \begin{scope}[shift={(56,0)}]
    % outside-the-image cells: left column and bottom row (zero)
    \foreach \r in {4,...,9} {
      \draw[ln-muted, densely dotted, fill=white] ({-1*\s},{-\r*\s}) rectangle ({0},{-(\r-1)*\s});
      \node[text=ln-muted] at ({-0.5*\s},{-(\r-0.5)*\s}) {0};
    }
    \foreach \c in {1,...,6} {
      \draw[ln-muted, densely dotted, fill=white] ({(\c-1)*\s},{-9*\s}) rectangle ({\c*\s},{-8*\s});
      \node[text=ln-muted] at ({(\c-0.5)*\s},{-8.5*\s}) {0};
    }
    % halo cells inside the image: top row from block (0,0), right column from (1,1)
    \foreach \c in {1,...,5} {
      \fill[gray!45, draw=white, line width=0.3pt] ({(\c-1)*\s},{-4*\s}) rectangle ({\c*\s},{-3*\s});
    }
    \fill[gray!45, draw=white, line width=0.3pt] ({5*\s},{-4*\s}) rectangle ({6*\s},{-3*\s});
    \foreach \r in {5,...,8} {
      \fill[gray!45, draw=white, line width=0.3pt] ({5*\s},{-\r*\s}) rectangle ({6*\s},{-(\r-1)*\s});
      \foreach \c in {1,...,5} {
        \fill[blue!30, draw=white, line width=0.3pt] ({(\c-1)*\s},{-\r*\s}) rectangle ({\c*\s},{-(\r-1)*\s});
      }
    }
    \draw[ln-muted] (0,-4*\s) rectangle (5*\s,-8*\s);
    \draw[ln-accent, thick] ({-1*\s},-3*\s) rectangle (6*\s,-9*\s);
    \foreach \c/\r/\t in {5/5/C, 5/4/1, 4/5/2, 6/5/3, 5/6/4} {
      \node[font=\scriptsize\sffamily\bfseries] at ({(\c-0.5)*\s},{-(\r-0.5)*\s}) {\t};
    }
    \node[anchor=south, text=ln-accent] at (2.5*\s,-3*\s+0.5)
      {\iflangja{共有メモリのタイル（ハローを含む）}{shared-memory tile with halo}};
    \node[anchor=west, text=ln-muted, align=left] at (6.5*\s,-5.5*\s)
      {\iflangja{灰色：他ブロックの要素，\\グローバルメモリから読む\\点線：画像の外，0 を入れる}{grey: other blocks' cells,\\read from global memory\\dotted: outside the image, zero}};
  \end{scope}
\end{tikzpicture}
   (body width, 120 mm)
\begin{tikzpicture}[x=1mm, y=1mm, font=\scriptsize\sffamily,
                    >={Stealth[length=1.4mm]}]
  \def\s{3.4}  % cell size
  % ---- left: image split among four blocks (rows counted from the top)
  \foreach \r in {1,...,8} {
    \foreach \c in {1,...,10} {
      \ifnum\r<5
        \ifnum\c<6 \def\col{orange!45}\else \def\col{yellow!60}\fi
      \else
        \ifnum\c<6 \def\col{blue!30}\else \def\col{green!45}\fi
      \fi
      \fill[\col, draw=white, line width=0.3pt]
        ({(\c-1)*\s},{-\r*\s}) rectangle ({\c*\s},{-(\r-1)*\s});
    }
  }
  \draw[ln-muted] (0,0) rectangle (10*\s,-8*\s);
  \draw[ln-muted, thick] (5*\s,0) -- (5*\s,-8*\s);
  \draw[ln-muted, thick] (0,-4*\s) -- (10*\s,-4*\s);
  \foreach \c/\r/\t in {5/5/C, 5/4/1, 4/5/2, 6/5/3, 5/6/4} {
    \node[font=\scriptsize\sffamily\bfseries] at ({(\c-0.5)*\s},{-(\r-0.5)*\s}) {\t};
  }
  \node[text=ln-muted] at (2.5*\s,1.8) {\iflangja{ブロック}{block} (0,0)};
  \node[text=ln-muted] at (7.5*\s,1.8) {\iflangja{ブロック}{block} (1,0)};
  \node[anchor=north, text=ln-muted] at (2.5*\s,-8*\s-0.5) {\iflangja{ブロック}{block} (0,1)};
  \node[anchor=north, text=ln-muted] at (7.5*\s,-8*\s-0.5) {\iflangja{ブロック}{block} (1,1)};
  % ---- right: the tile of block (0,1) with its halo
  \begin{scope}[shift={(56,0)}]
    % outside-the-image cells: left column and bottom row (zero)
    \foreach \r in {4,...,9} {
      \draw[ln-muted, densely dotted, fill=white] ({-1*\s},{-\r*\s}) rectangle ({0},{-(\r-1)*\s});
      \node[text=ln-muted] at ({-0.5*\s},{-(\r-0.5)*\s}) {0};
    }
    \foreach \c in {1,...,6} {
      \draw[ln-muted, densely dotted, fill=white] ({(\c-1)*\s},{-9*\s}) rectangle ({\c*\s},{-8*\s});
      \node[text=ln-muted] at ({(\c-0.5)*\s},{-8.5*\s}) {0};
    }
    % halo cells inside the image: top row from block (0,0), right column from (1,1)
    \foreach \c in {1,...,5} {
      \fill[gray!45, draw=white, line width=0.3pt] ({(\c-1)*\s},{-4*\s}) rectangle ({\c*\s},{-3*\s});
    }
    \fill[gray!45, draw=white, line width=0.3pt] ({5*\s},{-4*\s}) rectangle ({6*\s},{-3*\s});
    \foreach \r in {5,...,8} {
      \fill[gray!45, draw=white, line width=0.3pt] ({5*\s},{-\r*\s}) rectangle ({6*\s},{-(\r-1)*\s});
      \foreach \c in {1,...,5} {
        \fill[blue!30, draw=white, line width=0.3pt] ({(\c-1)*\s},{-\r*\s}) rectangle ({\c*\s},{-(\r-1)*\s});
      }
    }
    \draw[ln-muted] (0,-4*\s) rectangle (5*\s,-8*\s);
    \draw[ln-accent, thick] ({-1*\s},-3*\s) rectangle (6*\s,-9*\s);
    \foreach \c/\r/\t in {5/5/C, 5/4/1, 4/5/2, 6/5/3, 5/6/4} {
      \node[font=\scriptsize\sffamily\bfseries] at ({(\c-0.5)*\s},{-(\r-0.5)*\s}) {\t};
    }
    \node[anchor=south, text=ln-accent] at (2.5*\s,-3*\s+0.5)
      {\iflangja{共有メモリのタイル（ハローを含む）}{shared-memory tile with halo}};
    \node[anchor=west, text=ln-muted, align=left] at (6.5*\s,-5.5*\s)
      {\iflangja{灰色：他ブロックの要素，\\グローバルメモリから読む\\点線：画像の外，0 を入れる}{grey: other blocks' cells,\\read from global memory\\dotted: outside the image, zero}};
  \end{scope}
\end{tikzpicture}

  \caption{Left: a $10\times8$ image split among four blocks; $C$ and its
    neighbours 1--4 lie in three of them.  Right: the shared-memory tile
    of block (0,1): its own $5\times4$ region plus a one-cell halo.  Halo
    cells inside the image (grey) are loaded from global memory; those
    outside it (dotted) receive zero.}
  \label{fig:l03-stencil-tile}
\end{figure}

\begin{definition}[Boundary condition, zero padding]
  A \term{boundary condition}[境界条件] states what a stencil reads where
  a neighbour lies outside the array.  The simplest is
  \term{zero padding}[ゼロパディング]: positions outside the array are
  treated as zero, and
  the load phase stores \texttt{0} into the tile instead of reading
  global memory.  Other choices (clamping to the edge value, wrapping
  around, a one-sided formula) only change that branch of the load.
\end{definition}

\subsection{From stencil to convolution}

Rather than writing the average with four reads and an \texttt{if} for
each boundary, the kernel below is written as a small convolution: a
$3\times3$ filter is applied around every element, and the stencil is
expressed as the filter's coefficients.  This removes special cases and
makes the same kernel serve any $3\times3$ stencil.

\begin{definition}[Convolution, filter]
  A two-dimensional \term{convolution}[畳み込み] computes each output as a
  weighted sum of the input elements around the same position; the array
  of weights is the \term{filter}[フィルタ] (or kernel, a word avoided here
  because of its CUDA meaning).  Convolution in the signal-processing
  sense flips the filter; with a symmetric filter the two agree, and the
  machine-learning usage followed here does not flip.  With a filter $f$
  of radius $r$,
  \begin{equation}
    \texttt{out}[y][x] = \sum_{i=-r}^{r} \sum_{j=-r}^{r}
      \texttt{in}[y+i][x+j] \cdot f[i+r][j+r].
    \label{eq:l03-convolution}
  \end{equation}
  The four-neighbour average is the $3\times3$ filter with $0.25$ at the
  four edge-centres and $0$ elsewhere; a $3\times3$ filter is a 9-point
  stencil, and the halo it needs is $r = 1$ wide.
\end{definition}

\subsection{The kernel}

\begin{marginfigure}
  \resizebox{\linewidth}{!}{% Thread roles in one stencil block of 8 x 8 threads: the outer ring
% (grey) only loads halo cells into shared memory; the inner 6 x 6 threads
% (blue) load and compute one output each.
% Original (illustrates the code of the chapter, drawn at 8 x 8 instead
% of 16 x 16 for legibility).
% Usage: % Thread roles in one stencil block of 8 x 8 threads: the outer ring
% (grey) only loads halo cells into shared memory; the inner 6 x 6 threads
% (blue) load and compute one output each.
% Original (illustrates the code of the chapter, drawn at 8 x 8 instead
% of 16 x 16 for legibility).
% Usage: % Thread roles in one stencil block of 8 x 8 threads: the outer ring
% (grey) only loads halo cells into shared memory; the inner 6 x 6 threads
% (blue) load and compute one output each.
% Original (illustrates the code of the chapter, drawn at 8 x 8 instead
% of 16 x 16 for legibility).
% Usage: \input{figures/l03-tile-threads}   (margin width, 50 mm)
\begin{tikzpicture}[x=1mm, y=1mm, font=\scriptsize\sffamily]
  \def\s{5}
  \foreach \r in {0,...,7} \foreach \c in {0,...,7} {
    \ifnum\r=0 \def\col{gray!40}\else\ifnum\r=7 \def\col{gray!40}\else
      \ifnum\c=0 \def\col{gray!40}\else\ifnum\c=7 \def\col{gray!40}\else
        \def\col{blue!30}\fi\fi\fi\fi
    \fill[\col, draw=white, line width=0.4pt] (\c*\s,-\r*\s) rectangle ++(\s,-\s);
  }
  \draw[ln-accent, thick] (0,0) rectangle (8*\s,-8*\s);
  \draw[ln-muted] (\s,-\s) rectangle (7*\s,-7*\s);
  \node[anchor=south, text=ln-accent] at (4*\s,0.5)
    {\iflangja{ブロック：$8\times8$ スレッド}{block: $8\times8$ threads}};
  \node[anchor=north, text=ln-muted, align=center] at (4*\s,-8*\s-0.5)
    {\iflangja{灰色：読み込みのみ（ハロー）\\青：読み込み＋計算（$6\times6$ 出力）}{grey: load only (halo)\\blue: load and compute ($6\times6$ outputs)}};
\end{tikzpicture}
   (margin width, 50 mm)
\begin{tikzpicture}[x=1mm, y=1mm, font=\scriptsize\sffamily]
  \def\s{5}
  \foreach \r in {0,...,7} \foreach \c in {0,...,7} {
    \ifnum\r=0 \def\col{gray!40}\else\ifnum\r=7 \def\col{gray!40}\else
      \ifnum\c=0 \def\col{gray!40}\else\ifnum\c=7 \def\col{gray!40}\else
        \def\col{blue!30}\fi\fi\fi\fi
    \fill[\col, draw=white, line width=0.4pt] (\c*\s,-\r*\s) rectangle ++(\s,-\s);
  }
  \draw[ln-accent, thick] (0,0) rectangle (8*\s,-8*\s);
  \draw[ln-muted] (\s,-\s) rectangle (7*\s,-7*\s);
  \node[anchor=south, text=ln-accent] at (4*\s,0.5)
    {\iflangja{ブロック：$8\times8$ スレッド}{block: $8\times8$ threads}};
  \node[anchor=north, text=ln-muted, align=center] at (4*\s,-8*\s-0.5)
    {\iflangja{灰色：読み込みのみ（ハロー）\\青：読み込み＋計算（$6\times6$ 出力）}{grey: load only (halo)\\blue: load and compute ($6\times6$ outputs)}};
\end{tikzpicture}
   (margin width, 50 mm)
\begin{tikzpicture}[x=1mm, y=1mm, font=\scriptsize\sffamily]
  \def\s{5}
  \foreach \r in {0,...,7} \foreach \c in {0,...,7} {
    \ifnum\r=0 \def\col{gray!40}\else\ifnum\r=7 \def\col{gray!40}\else
      \ifnum\c=0 \def\col{gray!40}\else\ifnum\c=7 \def\col{gray!40}\else
        \def\col{blue!30}\fi\fi\fi\fi
    \fill[\col, draw=white, line width=0.4pt] (\c*\s,-\r*\s) rectangle ++(\s,-\s);
  }
  \draw[ln-accent, thick] (0,0) rectangle (8*\s,-8*\s);
  \draw[ln-muted] (\s,-\s) rectangle (7*\s,-7*\s);
  \node[anchor=south, text=ln-accent] at (4*\s,0.5)
    {\iflangja{ブロック：$8\times8$ スレッド}{block: $8\times8$ threads}};
  \node[anchor=north, text=ln-muted, align=center] at (4*\s,-8*\s-0.5)
    {\iflangja{灰色：読み込みのみ（ハロー）\\青：読み込み＋計算（$6\times6$ 出力）}{grey: load only (halo)\\blue: load and compute ($6\times6$ outputs)}};
\end{tikzpicture}
}
  \caption{Thread roles in one block (drawn for $8\times8$ threads): the
    outer ring only loads halo cells; the inner threads load and compute.}
  \label{fig:l03-tile-threads}
\end{marginfigure}

The design below launches a block of $(\mathit{TILE}+2r)^2$ threads for
a tile of $\mathit{TILE}^2$ outputs: every thread loads one element of
the tile including the halo, and only the inner threads, those whose
element is not a halo cell, compute and store an output
(\cref{fig:l03-tile-threads}).  The filter lives in constant memory.

\begin{lstlisting}[language=C++, numbers=none]
#define TILE 14               // outputs per block side
#define RADIUS 1              // filter radius r
#define BLOCK (TILE + 2 * RADIUS) // 16 per side

__constant__ float filter[2 * RADIUS + 1][2 * RADIUS + 1];

__global__ void stencil2D(const float *in, float *out,
                          int width, int height) {
  __shared__ float tile[BLOCK][BLOCK];
  int tx = threadIdx.x, ty = threadIdx.y;
  int x = blockIdx.x * TILE + tx - RADIUS;  // image column
  int y = blockIdx.y * TILE + ty - RADIUS;  // image row

  // 1. every thread loads one cell, zero outside the image
  if (x >= 0 && x < width && y >= 0 && y < height)
    tile[ty][tx] = in[y * width + x];
  else
    tile[ty][tx] = 0.0f;
  __syncthreads();

  // 2. inner threads compute and 3. store one output each
  bool inner = tx >= RADIUS && tx < BLOCK - RADIUS &&
               ty >= RADIUS && ty < BLOCK - RADIUS;
  if (inner && x < width && y < height) {
    float acc = 0.0f;
    for (int i = -RADIUS; i <= RADIUS; ++i)
      for (int j = -RADIUS; j <= RADIUS; ++j)
        acc += tile[ty + i][tx + j]
             * filter[i + RADIUS][j + RADIUS];
    out[y * width + x] = acc;
  }
}
\end{lstlisting}
The host sets the filter once with \texttt{cudaMemcpyToSymbol} and
launches with a two-dimensional configuration:
\begin{lstlisting}[language=C++, numbers=none]
dim3 block(BLOCK, BLOCK);   // 16 x 16 threads
dim3 grid((width + TILE - 1) / TILE,
          (height + TILE - 1) / TILE);
stencil2D<<<grid, block>>>(d_in, d_out, width, height);
\end{lstlisting}
The index arithmetic is \cref{eq:l03-global-index} shifted by $r$: thread
$(0,0)$ of block $(0,0)$ loads image position $(-1,-1)$, which is outside
and becomes zero, and thread $(1,1)$ loads $(0,0)$, the first output.
The offset of one block is \texttt{TILE}, not \texttt{BLOCK}, because
neighbouring tiles overlap by the halo.\caution{One barrier suffices
here: the tile is loaded once and only read afterwards.  A kernel that
reuses the shared buffer for several tiles needs a second
\texttt{\_\_syncthreads()} after the compute phase.}

\begin{example}[A $64\times64$ image]
  \label{ex:l03-stencil-count}
  With $\mathit{TILE}=14$ and $r=1$ on a $64\times64$ image:
  \begin{itemize}
    \item grid $\lceil 64/14 \rceil \times \lceil 64/14 \rceil = 5\times5
      = 25$ blocks of $16\times16 = 256$ threads, \num{6400} threads for
      \num{4096} outputs; in each of the 16 interior blocks 196 threads
      compute and 60 only load;
    \item shared memory per block $16 \times 16 \times 4 =
      \SI{1}{\kibi\byte}$;
    \item tile loads: $25 \times 256 = \num{6400}$, of which
      $72^2 = \num{5184}$ read global memory (the rest are zero fills;
      per dimension the blocks read $15+16+16+16+9 = 72$ positions),
      against $9 \times \num{4096} = \num{36864}$ if every thread read its
      nine inputs from global memory, a factor of about~7; the halo
      overhead is $256/196 = 1.31$ loads per output for an interior block
      ($\num{5184}/\num{4096} \approx 1.27$ over the whole image).
  \end{itemize}
  The last block in each dimension covers columns 56--69 of a 64-wide
  image: its threads with $x \ge 64$ load zeros and skip the store.\tip{Halo fraction $1 - (\mathit{TILE}/\mathit{BLOCK})^2$: 23\,\% of
the threads here, 12\,\% for a $32\times32$ block.  Larger tiles waste
fewer threads but need more shared memory and leave fewer blocks per SM
(\cref{sec:l03-occupancy}).}
\end{example}

\Needspace{16\baselineskip}

The alternative design, in which the block has exactly $\mathit{TILE}^2$
threads and the halo is loaded by extra passes of the edge threads,
launches no idle threads but needs more conditional code; PMPP ch.~8
shows both.\margin{All threads of a warp read the same filter
coefficient at once; the constant cache serves one address to the whole
warp as a broadcast, which is what constant memory is for.}

\begin{keypoint}[Lesson summary]
  A CUDA program is host code that allocates device memory
  (\texttt{cudaMalloc}), copies data in and out (\texttt{cudaMemcpy}) and
  launches kernels with an execution configuration
  \texttt{<<<blocks, threads>>>}.  A kernel runs as a grid of
  independent thread blocks; each thread computes its global index from
  \texttt{blockIdx}, \texttt{blockDim} and \texttt{threadIdx}, blocks
  run in no fixed order, and the threads of one block share an on-chip
  shared memory and synchronise with \texttt{\_\_syncthreads()} in the
  BSP pattern load, barrier, compute, store.  The five memory spaces,
  local, shared, global, constant and texture, follow the thread
  hierarchy in scope.  The number of blocks resident on an SM is the
  minimum allowed by its thread slots, registers and shared memory, and
  occupancy is the fraction of the SM's warp slots they fill; higher
  occupancy hides latency better but is not always faster.  A stencil
  becomes a tiled kernel: each block loads its tile plus a halo into
  shared memory, with zero padding at the image edge, synchronises once,
  and computes one output per inner thread.
\end{keypoint}

\end{document}
